\begin{enumerate}
\def\labelenumi{\arabic{enumi}.}
\setcounter{enumi}{10}
\item
  设 \(\mathfrak{g}\) 为一个李代数，\(\mathfrak{h}\) 为 \(\mathfrak{g}\) 的理想，且 \(\mathcal{D}\mathfrak{h} = \mathfrak{h}\) 。证明 \(\mathfrak{h}\) 是 \(\mathfrak{g}\) 的特征理想。
\item
  设 g 为一个李代数。
\end{enumerate}

\begin{enumerate}
\def\labelenumi{(\alph{enumi})}
\item
  导子 D 与 g 的所有内导子交换的充要条件是 D 将 g 映入其中心。
\item
  以下设 \(\mathfrak{g}\) 的中心为零，于是将 \(\mathfrak{g}\) 视为由 \(\mathfrak{D}\) 的所有导子组成的李代数 \(\mathfrak{g}\) 的一个理想。证明 \(\mathfrak{g}\) 在 \(\mathfrak{D}\) 中的中心化子为零（使用 (a)）。特别地，\(\mathfrak{D}\) 的中心为零。
\item
  证明 \(\Delta\) 的满足 \(\mathfrak{D}\) 的导子 \(\Delta(\mathfrak{g}) = \{0\}\) 为零。 \(([\Delta(\mathfrak{D}),\mathfrak{g}]\subset \Delta([\mathfrak{D},\mathfrak{g}]) + [\mathfrak{D},\Delta(\mathfrak{g})]\subset \Delta(\mathfrak{g}) = \{0\}\) 因而由 (b) 得 \(\Delta(\mathfrak{D}) = \{0\}\)。)
\item
  进一步证明，若 \(\mathfrak{g} = \mathcal{D}\mathfrak{g}\)，则 \(\Delta\) 的每个导子 \(\mathfrak{D}\) 都是内导子（使用 (c) 和练习 11）。
\end{enumerate}

\begin{enumerate}
\def\labelenumi{\arabic{enumi}.}
\setcounter{enumi}{12}
\tightlist
\item
  设 \(\mathfrak{g}\) 为一个李代数，\(\mathfrak{a}\) 和 \(\mathfrak{b}\) 为 \(\mathfrak{g}\) 的两个子模。对于每个整数 \(i \geqslant 0\)，如下定义子模 \(\mathfrak{m}_i = \mathfrak{m}_i(\mathfrak{a},\mathfrak{b})\) 和 \(\mathfrak{n}_i = \mathfrak{n}_i(\mathfrak{a},\mathfrak{b})\)：\(\mathfrak{m}_1 = \mathfrak{b}, \mathfrak{m}_{i+1} = [\mathfrak{m}_i,\mathfrak{b}], \mathfrak{n}_1 = \mathfrak{a}, \mathfrak{n}_{i+1} = [\mathfrak{n}_i,\mathfrak{b}]\)。证明
\end{enumerate}

\[
[ a, m _ {i} ] \subset n _ {i + 1}
\]

对于 \(i = 1,2,\ldots\)（对 \(i\) 作归纳，注意到

\[
\mathfrak {m} _ {i} ([ a, b ], b) = \mathfrak {m} _ {i} (a, b)
\]

以及 \(n_{i}([a,b],b)=n_{i+1}(a,b)\)。

¶ 14. 设 \(\mathfrak{a}\) 为一个李代数，\(\mathfrak{b}\) 为 \(\mathfrak{a}\) 的子代数。连接 \(\mathfrak{a}\) 与 \(\mathfrak{b}\) 的合成列是一个递减序列 \((\mathfrak{a}_i)_{0\leqslant i\leqslant n}\)，其中各项是 \(\mathfrak{a}\) 的子代数，满足 \(\mathfrak{a}_0 = \mathfrak{a}\)、\(\mathfrak{a}_n = \mathfrak{b}\)，且 \(\mathfrak{a}_{i+1}\) 是 \(\mathfrak{a}_i\) 的理想，其中 \(0\leqslant i < n\)。称 \(\mathfrak{b}\) 为 \(\mathfrak{a}\) 的次不变子代数，当且仅当存在连接 \(\mathfrak{a}\) 与 \(\mathfrak{b}\) 的合成列。

\begin{enumerate}
\def\labelenumi{(\alph{enumi})}
\item
  设 \(\mathfrak{b}\) 是 \(\mathfrak{a}\) 的次不变子代数，\((\mathfrak{a}_i)_{0\leqslant i\leqslant n}\) 是连接 \(\mathfrak{a}\) 与 \(\mathfrak{b}\) 的合成列。由练习 13 推出 \([\mathcal{C}^{n+p}\mathfrak{b},\mathfrak{a}]\subset\mathcal{C}^{p}\mathfrak{b}\)，注意到当 \([\mathfrak{a}_i,\mathfrak{b}]\subset\mathfrak{a}_{i+1}\) 时 \(1\leqslant i<n\)。
\item
  由 (a) 推出，所有 \(\mathcal{C}^{\infty}\mathfrak{b}\)（\(\mathcal{C}^{p}\mathfrak{b}\)）的交 \(p = 0,1,2,\ldots\) 是 \(\mathfrak{a}\) 的理想。
\item
由 (b) 推出，若 \(\mathfrak{c}\) 是 \(\mathfrak{a}\) 的次不变子代数且满足 \(\mathfrak{c} = \mathcal{D}\mathfrak{c}\)，则它是 \(\mathfrak{a}\) 的理想。
\item
  当 K 是一个域且 \(\dim_{K}a<+\infty\) 时，证明所有 \(D^{\infty}b\) 的交 \(D^{p}b\left(p=0,1,2,\ldots\right)\) 是 a 的理想。（证明此交是 a 的次不变子代数，并应用 (c)。）
\end{enumerate}

¶ 15. (a) 设 a 为一个李代数，b 为 a 的子代数，c 为 b 的理想，z 为 c 在 a 中的中心化子。则 {[}b, b{]} ⊂ z。推出 b + z 是 a 的子代数，且 z 是其中的理想。

\begin{enumerate}
\def\labelenumi{(\alph{enumi})}
\setcounter{enumi}{1}
\item
  设 g 为一个李代数，a 为 g 的次不变子代数，b 为 a 的理想。若 b 在 a 中的中心化子以及 a 在 g 中的中心化子均为零，则 \(\mathfrak{z}\)（\(\mathfrak{b}\) 在 \(\mathfrak{g}\) 中的中心化子）为零。（令 \(\mathfrak{b} = \mathfrak{a} + \mathfrak{z}\)，由 \((a)\) 可知它是子代数；\(\mathfrak{a}\) 在 \(\mathfrak{b}\) 中次不变；若 \(\mathfrak{z} \neq \{0\}\)，则 \(\mathfrak{a} \neq \mathfrak{h}\)，因而 \(\mathfrak{a}\) 是 \(\mathfrak{a}'\) 的理想，其中 \(\mathfrak{a} \neq \mathfrak{a}' \subset \mathfrak{b}\)；取 \(a' \in \mathfrak{a}'\)、\(a' \notin \mathfrak{a}\)，并令 \(a' = a + z\)，其中 \(a \in \mathfrak{a}\)、\(z \in \mathfrak{z}\)、\(z \neq 0\)；证明 \([z, \mathfrak{a}]\) 包含于 \(\mathfrak{a}\) 且与 \(\mathfrak{b}\) 交换，于是 \([z, \mathfrak{a}] = \{0\}\)，矛盾。)
\item
设 \(\mathfrak{g}\) 为中心为零的李代数，\(\mathfrak{D}_1\) 为 \(\mathfrak{g}\) 的导子李代数，\(\mathfrak{D}_2\) 是 \(\mathfrak{D}_1, \ldots, \mathfrak{g}\) 的导子李代数，后者被识别为 \(\mathfrak{D}_1\) 的理想，\(\mathfrak{D}_1\) 被识别为 \(\mathfrak{D}_2, \ldots\) 的理想，……（参见练习 12 (b)）。利用 (b) 和练习 12 (b)，证明 \(\mathfrak{g}\) 在 \(\mathfrak{D}_i\) 中的中心化子对所有 \(i\) 均为零。（关于本练习的后续内容，参见 § 4，练习 15。）
\end{enumerate}

\begin{enumerate}
\def\labelenumi{\arabic{enumi}.}
\setcounter{enumi}{15}
\tightlist
\item
  设 g 为一个李代数，D 为 g 的导子李代数。D 的恒等映射定义了 D 与 g 的半直积 h，称为 g 的全纯代数。将 g 视为 h 的一个理想，将 D 视为 h 的一个子代数。
\end{enumerate}

\begin{enumerate}
\def\labelenumi{(\alph{enumi})}
\item
  证明 \(\mathfrak{g}^*\) 在 \(\mathfrak{g}\) 中的中心化子 \(\mathfrak{h}\) 是集合 \(\operatorname{ad}_{\mathfrak{g}}x - x(x\in \mathfrak{g})\)。
\item
  对于 \(x + d\) 的每个元素 \(\mathfrak{h}\)（\(x \in \mathfrak{g}, d \in \mathfrak{D}\)），记
\end{enumerate}

\[
\theta (x + d) = \mathrm{ad} _ {\mathfrak {g}} x - x + d.
\]

证明 \(\theta\) 是 \(\mathfrak{h}\) 的二阶自同构，并将 \(\mathfrak{g}\) 映到 \(\mathfrak{g}^*\)，因而可以将 \(\mathfrak{h}\) 看作 \(\mathfrak{g}^*\) 的全纯代数。

\begin{enumerate}
\def\labelenumi{(\alph{enumi})}
\setcounter{enumi}{2}
\item
  对于每个理想 \(\mathfrak{k} \neq \{0\}\)，若其为 \(\mathfrak{h}\) 的理想，则 \(\mathfrak{k} \cap (\mathfrak{g} + \mathfrak{g}^*) \neq \{0\}\)。（若 \(\mathfrak{k} \cap \mathfrak{g} = \{0\}\)，则 \(\mathfrak{k} \subset \mathfrak{g}^*\)。）
\item
  若 g 是交换的，则 h 的每个非零理想 \(t \neq \{0\}\) 都包含 g。若 K 是一个域且 \(\dim_{K} g < +\infty\)，推出当 \(t \neq g\) 时，h 不可能是 t 的全纯代数（使用 (b)）。
\end{enumerate}

\begin{enumerate}
\def\labelenumi{\arabic{enumi}.}
\setcounter{enumi}{16}
\tightlist
\item
  设 K 是一个域。令 T 为一个不定元，且 \(L = K((T))\)；L 具有一个规范赋值（形式幂级数的阶），从而成为完备赋值域。
\end{enumerate}

\begin{enumerate}
\def\labelenumi{(\alph{enumi})}
\item
  设 \(\mathrm{V}\) 是 \(\mathbf{K}\) 上的有限维向量空间。则 \(\mathrm{V}_{(\mathbb{L})}\) 规范地具有完备可度量化的拓扑向量空间结构（《拓扑向量空间》，第 I 章第 2 节第 3 号，定理 2）。设 \((x_{p})_{p\in \mathbb{Z}}\) 是 \(\mathrm{V}\) 中元素组成的族，且 \(x_{p} = 0\) 对于小于某个整数的 \(p\) 成立。则级数 \(\sum_{-\infty}^{+\infty}x_p\mathrm{T}^p\) 在 \(\mathrm{V}_{(\mathrm{L})}\) 中收敛，并且 \(\mathrm{V}_{(\mathrm{L})}\) 的每个元素都能唯一地以这种方式表示。（取 V 的一个基。）
\item
  向量空间 \(\mathcal{L}_{\mathrm{L}}(\mathrm{V}_{(\mathrm{L})}) = \mathcal{L}_{\mathrm{K}}(\mathrm{V})_{(\mathrm{L})}\) 具有规范的完备可度量化拓扑向量空间结构。 \(\mathrm{V}_{(\mathrm{L})}\) 的每个自同态都能唯一地表示为 \(\sum_{-\infty}^{+\infty}u_p\mathrm{T}^p\) 的形式，其中 \(u_{p}\in \mathcal{L}_{\mathrm{K}}(\mathrm{V})\)，且当 \(u_{p}\) 小于某个整数时 \(p\) 为零。
\item
  如果 \(\mathbf{K}\) 的特征为 0 且 \(u\in \mathcal{L}_{\mathrm{K}}(\mathrm{V})\)，则记作
\end{enumerate}

\[
e ^ {\mathrm{Tu}} = 1 + \frac {1}{1 !} u \mathrm{T} + \frac {1}{2 !} u ^ {2} \mathrm{T} ^ {2} + \dots \in \mathcal {L} _ {\mathrm{L}} (\mathrm{V} _ {(\mathrm{L})}).
\]

若 \(x \in V\) 满足 \(ux = 0\)，则 \(e^{Tu}.x = x\)。

\begin{enumerate}
\def\labelenumi{(\alph{enumi})}
\setcounter{enumi}{3}
\item
  若 W 是 K 上另一个有限维向量空间，且 \(v \in \mathcal{L}_{\mathrm{K}}(\mathrm{W})\)，则 \(e^{\mathrm{T}(u \otimes 1 + 1 \otimes v)} = e^{\mathrm{T}u} \otimes e^{\mathrm{T}v}\)。
\item
  由 (c) 和 (d) 推出，若 V 具有一个（不必结合的）代数结构且 u 是 V 的导子，则 \(e^{Tu}\) 是代数 \(V_{(L)}\) 的自同构。
\end{enumerate}

¶ 18. 设 H 是一个复 Hilbert 空间，Y 和 Z 是 H 上的连续厄米算子，且 B = {[}Z, Y{]}。设 Y 与 B 可交换。

\begin{enumerate}
\def\labelenumi{(\alph{enumi})}
\item
  证明 \([Z, Y^n] = nBY^{n-1}\)（对 \(n\) 作归纳）。
\item
  设 \(f\) 是实变量的连续可微函数。由 (a) 推出 \([Z, f(Y)] = Bf'(Y)\)。
\item
  由 (b) 推出 \(B = 0\)。（证明若 \(B \neq 0\)，可以选取 \(f\) 使得 \(\| f(Y) \| \leqslant 1\)，而 \(\| Bf'(Y) \|\) 任意大。）
\item
  设 g 是 H 上满足 \(T^{*} = -T\) 的连续算子 T 所成的李代数。由 (c) 推出，条件 \(Y \in g\)、\(Z \in g\)、\([[Z, Y], Y] = 0\) 蕴含 \([Z, Y] = 0\)。
\end{enumerate}

\begin{enumerate}
\def\labelenumi{\arabic{enumi}.}
\setcounter{enumi}{18}
\tightlist
\item
  \begin{enumerate}
  \def\labelenumii{(\alph{enumii})}
  \tightlist
  \item
    设 \(\mathbf{L}\) 是 \(\mathbf{K}\) 上的结合代数。给定 \(k\) 个元素 \(x_{1},\ldots ,x_{k}\) 属于 \(\mathbf{L}\)，记：
  \end{enumerate}
\end{enumerate}

\[
f _ {k} (x _ {1}, \dots , x _ {k}) = \sum_ {\sigma \in \mathfrak {S}_{k}} x _ {\sigma (1)} \dots x _ {\sigma (k)}.
\]

在代数 \(L \otimes_{K} K[T_{1}, \ldots, T_{k}]\) 中，其中 \(T_{i}\) 是不定元，\(f_{k}(x_{1}, \ldots, x_{k})\) 是 \(T_{1} \ldots T_{k}\) 中 \((x_{1}T_{1} + \cdots + x_{k}T_{k})^{k}\) 的系数；它也是下式中 \(T_{1} \ldots T_{k}\) 的系数：

\[
\sum_ {j = 0} ^ {k - 1} \left(x _ {1} \mathrm{T} _ {1} + \dots + x _ {k - 1} \mathrm{T} _ {k - 1}\right) ^ {j} x _ {k} \mathrm{T} _ {k} \left(x _ {1} \mathrm{T} _ {1} + \dots + x _ {k - 1} \mathrm{T} _ {k - 1}\right) ^ {k - j - 1}.
\]

  给定 \(x, y\)（它们是 \(\mathbf{L}\) 中的两个元素），定义 \(g_{i,k-i}(x,y)\) 为 \(\mathrm{T}_1^i\mathrm{T}_2^{k - i}\) 中 \((x\mathrm{T}_1 + y\mathrm{T}_2)^k\) 的系数。证明

\[
i! (k - i)! g _ {i, k - i} (x, y) = f _ {k} \left(t _ {1}, \dots , t _ {k}\right)
\]

其中当 \(t_h = x\) 时 \(h \leqslant i\)，当 \(t_h = y\) 时 \(h \geqslant i + 1\)。

\begin{enumerate}
\def\labelenumi{(\alph{enumi})}
\setcounter{enumi}{1}
\tightlist
\item
  若将 L 视为 K 上的李代数，则在 \(\mathcal{L}_{\mathrm{K}}(\mathrm{L})\) 中：
\end{enumerate}

\[
(\operatorname{ad} x) ^ {n} = \sum_ {i = 0} ^ {n} (- 1) ^ {n - i} \binom {n} {i} \mathrm{L} _ {x} ^ {i} R _ {x} ^ {n - i},
\]

其中 \(\mathbf{L}_x\)（分别为 \(R_{x}\)）是乘法 \(y\mapsto xy\)（分别为 \(y\mapsto yx\)）。

推出，若 K 的特征为 p（p 为素数），则：

\begin{enumerate}
\def\labelenumi{(\arabic{enumi})}
\tightlist
\item
\end{enumerate}

\[
(\operatorname{ad} x) ^ {p}. y = (\operatorname{ad} x ^ {p}). y\tag{2}
\]

\[
(\operatorname{ad} x) ^ {p - 1}. y = \sum_ {i = 0} ^ {p - 1} x ^ {i} y x ^ {p - 1 - i}.
\]

由 (2) 推出：

\[
f _ {p - 1} (\operatorname{ad} x _ {1}, \dots , \operatorname{ad} x _ {p - 1}). y = f _ {p} (x _ {1}, \dots , x _ {p - 1}, y)
\]

以及：

\[
g _ {i, p - 1 - i} (\mathrm{ad} x, \mathrm{ad} y). y = (p - i) g _ {i, p - i} (x, y).
\]

从而对 L 中任意两个元素 x、y：

\[
(x + y) ^ {p} = x ^ {p} + y ^ {p} + \Lambda_ {p} (x, y)\tag{3}
\]

其中：

\[
\Lambda_ {p} (x, y) = \sum_ {i = 1} ^ {p - 1} (p - i) ^ {- 1} g _ {i, p - 1 - i} (\mathrm{ad} x, \mathrm{ad} y). y
\]

属于由 x 和 y 生成的 L 的李子代数（“Jacobson 公式”）。

\begin{enumerate}
\def\labelenumi{\arabic{enumi}.}
\setcounter{enumi}{19}
\tightlist
\item
  设 g 是环 K 上满足 \(pK = \{0\}\) 的李代数（p 为素数）。g 到自身的映射 \(x \mapsto x^{[p]}\) 称为 p-映射，如果它满足关系式
\end{enumerate}

\[
\begin{array}{r l} & {\operatorname{ad} x ^ {[ p ]} = (\operatorname{ad} x) ^ {p}} \\ & {(\lambda x) ^ {[ p ]} = \lambda^ {p} x ^ {[ p ]}} \\ & {(x + y) ^ {[ p ]} = x ^ {[ p ]} + y ^ {[ p ]} + \Lambda_ {p} (x, y) \quad (\mathrm{cf.Exercise19(b)})} \end{array}
\]

对于 \(x, y\)（属于 \(\mathfrak{g}\)）以及 \(\lambda \in \mathbf{K}\)。李 \(p\)-代数定义在 \(\mathbf{K}\) 上，是赋予其一个 \(\mathbf{K}\) 上的李代数结构和一个 \(p\)-映射所得到的结构。每个 \(\mathbf{K}\) 上的结合代数都是李 \(p\)-代数，其中 \([x, y] = xy - yx\) 且 \(x^{[p]} = x^p\)（练习 19）。映射 \(u\) 将李 \(p\)-代数 \(\mathfrak{g}\) 映入李 \(p\)-代数 \(\mathfrak{g}'\)；称为 \(p\)-同态，只要 \(u\) 是李代数同态且 \(u(x^{[p]}) = (u(x))^{[p]}\)。证明，在李 \(p\)-代数 \(\mathfrak{g}\) 中，每个 \(p\)-映射都具有 \(x \mapsto x^{[p]} + f(x)\) 的形式，其中 \(f\) 是将 \(\mathfrak{g}\) 映入其中心的映射，并且关于自同态 \(\lambda \mapsto \lambda^p\)（作用于 \(\mathbf{K}\)）是半线性的。更一般地，若 \(u\) 是从 \(\mathfrak{g}\) 到李 \(p\)-代数的同态，其中 \(\mathfrak{g}', u(x^{[p]}) - (u(x))^{[p]}\) 属于 \(\mathfrak{g}'\) 的中心。

\begin{enumerate}
\def\labelenumi{\arabic{enumi}.}
\setcounter{enumi}{20}
\tightlist
\item
  \begin{enumerate}
  \def\labelenumii{(\alph{enumii})}
  \tightlist
  \item
    设 \(\mathbf{K}\) 是特征为 \(p > 0\) 的域，\(\mathbf{L}\) 是 \(\mathbf{K}\) 上不必结合的代数。证明 \(\mathbf{L}\) 的导子构成一个李 \(p\)-代数，其 \(p\)-映射为 \(\mathbf{D} \mapsto \mathbf{D}^p\)。
  \end{enumerate}
\end{enumerate}

\begin{enumerate}
\def\labelenumi{(\alph{enumi})}
\setcounter{enumi}{1}
\item
  设 \(\mathbf{L}\) 为代数 \(\mathrm{K}[\mathrm{X}] / (\mathrm{X}^p)\)，\(\mathfrak{w}\) 为 \(p\) 的导子所成的 \(\mathbf{L}\)-代数。令 \(\mathrm{D}_i\)（\(0 \leqslant i \leqslant p - 1\)）表示满足 \(\mathbf{L}\) 的 \(\mathrm{D}_i(\mathbf{X}) = \mathbf{X}^i\) 的导子。证明 \(\mathrm{D}_i\) 构成 \(\mathfrak{w}\) 在 \(\mathbf{K}\) 上的一组基，并且当 \([\mathrm{D}_i, \mathrm{D}_j] = (j - i)\mathrm{D}_{i + j - 1}\) 时 \(i + j \leqslant p\)，当 \([\mathrm{D}_i, \mathrm{D}_j] = 0\) 时 \(i + j > p\)，且 \(\mathrm{D}_i^p = 0\)；但 \(i = 1\) 时例外，此时 \(\mathrm{D}_1^p = \mathrm{D}_1\)。
\item
  证明，若 \(p \geqslant 3\)，李代数 \(\mathfrak{w}\) 只有理想 \(\{0\}\) 和 \(\mathfrak{w}\)。（利用 \(D_i\) 的乘法表，证明 \(\mathfrak{w}\) 的每个非零理想都包含 \(D_{p-1}\) 的一个非零倍数。）
\item
令 \(V_{k}\) 为 w 的子空间，其基由 \(D_{i}\)（指标为 \(\geqslant k\)）构成。证明，若 \(p \geqslant 5\)，则 \(V_k\) 恰为 \(Z \in w\) 中中心化子维数为 \(\geqslant k + 1\) 的元素所成的集合。推出，若 \(p \geqslant 5\)，\(w\) 的自同构群是可解的。
\end{enumerate}

\begin{enumerate}
\def\labelenumi{\arabic{enumi}.}
\setcounter{enumi}{21}
\tightlist
\item
  设 g 是特征为 p 的环 K 上的李 p-群（p 为素数）。令 \(x \mapsto x^{[p]}\) 表示 p-映射。对于 g 的每个子集 E，令 \(E^{p}\) 表示由 \(x^{p}\)（其中 \(x \in E\)）生成的 g 的子模。理想 \(a \subset g\) 若满足 \(a^{p} \subset a\)，则称为 p-理想。
\end{enumerate}

\begin{enumerate}
\def\labelenumi{(\alph{enumi})}
\item
对于理想 \(\mathfrak{a} \subset \mathfrak{g}\)，它为 \(p\)-理想的充要条件是它是某个 \(p\)-同态的核（练习 20）。
\item
两个 \(p\)-理想之和是 \(p\)-理想。\(p\)-理想与 \(p\)-子代数之和是 \(p\)-子代数。
\item
设 \(\mathfrak{h}\) 是 \(\mathfrak{g}\) 的李子代数。最小李 \(p\)-子代数（包含 \(\mathfrak{h}\)）是 \(\mathfrak{h} + \mathfrak{h}^p + \mathfrak{h}^{p^2} + \cdots = \overline{\mathfrak{h}}\)；若记 \(\mathfrak{h}_i = \mathfrak{h} + \mathfrak{h}^p + \cdots + \mathfrak{h}^{p^i}\)，则 \(\mathfrak{h}_i\) 是 \(\overline{\mathfrak{h}}\) 的理想，且 \(\overline{\mathfrak{h}}/\mathfrak{h}\) 是交换的。若 \(\mathfrak{h}\) 是 \(\mathfrak{g}\) 的理想，则 \(\mathfrak{h}_i\) 和 \(\overline{\mathfrak{h}}\) 也是理想，并且后者是最小 \(p\)-理想，包含 \(\mathfrak{h}\)。
\item
  g 的任意子模的正规化子和中心化子都是 g 的 p-子代数。
\end{enumerate}

\begin{enumerate}
\def\labelenumi{\arabic{enumi}.}
\setcounter{enumi}{22}
\tightlist
\item
设 \(\mathfrak{g}\) 是一个有限维交换李 \(p\)-代数，定义在完美域 \(\mathbf{K}\) 上，且其特征为 \(p > 0\)。
\end{enumerate}

\begin{enumerate}
\def\labelenumi{(\alph{enumi})}
\item
  证明 g 唯一地是两个 p-子代数 h、t 的直和，其中 p-映射在 h 上是双射，在 t 上是幂零的（考察 g 上 p-映射的迭代；参见《代数学》，第 VIII 章第 2 节第 2 号引理和第 VII 章第 5 节练习 20 (a)）。h 称为 g 的 p-核。
\item
  设 K 是代数闭域。证明存在一组基 \((e_{i})\)，使 \(\mathfrak{h}\) 中的这些基元满足 \(e_{i}^{p}=e_{i}\)（考察由单个元素生成的 h 的 p-子代数，证明它包含某个满足 \(x^{p}=x\neq0\) 的 x；然后对 \(\mathfrak{h}\) 的维数作归纳）。推出，h 的 p-子代数只有有限多个，并且它们与由 \(F_{p}\) 生成的素域 \(e_{i}\) 上的向量空间的向量子空间一一对应。
\item
证明存在一组基 \((f_{ij})\)，属于 \(\mathfrak{t}\)（\(1 \leqslant i \leqslant k\)、\(1 \leqslant j \leqslant s_i\)，对所有 \(i\)），其中 \(s_i\) 的序列递减，并满足 \(f_{1j}^p = 0\)（当 \(1 \leqslant j \leqslant s_1\) 时），\(f_{ij}^p = f_{i-1,j}\)（当 \(2 \leqslant i \leqslant k\)、\(1 \leqslant j \leqslant s_i\) 时）（方法见《代数学》，第 VII 章第 5 节练习 20 (b)）。
\end{enumerate}

¶ 24. 设 K 为一个（交换）域，具有（任意）特征 p，并设 \(X_{i}, Y_{i}, Z_{i}\) 是 3n 个彼此不同的不定元；为简便起见，分别将系统 \((\mathrm{X}_{i})_{1 \leq i \leq n}, (\mathrm{Y}_{i})_{1 \leq i \leq n}, (\mathrm{Z}_{i})_{1 \leq i \leq n}\) 记为 x、y、z。令 \(A_{x,y,z}\) 为系数在 K 中、关于 3n 个不定元 \(X_{i}, Y_{i}, Z_{i}\) 的形式幂级数代数；将 \(A_{x,y}\)（分别为 \(A_{x}\)）记作 \(A_{x,y,z}\) 的不含 z（分别不含 y、z）的幂级数所成的子代数。\(A_{x,y,z}\)（分别为 \(A_{x,y}, A_{x}\)）中的元素记为 \(u(\mathbf{x}, \mathbf{y}, \mathbf{z})\)（分别记为 \(u(\mathbf{x}, \mathbf{y}), u(\mathbf{x})\)）；

  一组 \((u_{i}(\mathbf{x},\mathbf{y},\mathbf{z}))_{1\leqslant i\leqslant n}\) 由 \(A_{x,y,z}\) 中的 n 个元素组成，记为 \(u(\mathbf{x},\mathbf{y},\mathbf{z})\)；对于只含三个系统 x、y、z 中一个或两个的形式幂级数，也采用类似记号。如果 \(u(\mathbf{x},\mathbf{y},\mathbf{z})\in A_{x,y,z}\)，且 \(\mathbf{f}=(f_{i}),\mathbf{g}=(g_{i}),\mathbf{h}=(h_{i})\) 是由 \(A_{x,y,z}\) 中没有常数项的 n 个形式幂级数组成的三个系统，则用 \(u(\mathbf{f}(\mathbf{x},\mathbf{y},\mathbf{z}),\mathbf{g}(\mathbf{x},\mathbf{y},\mathbf{z}),\mathbf{h}(\mathbf{x},\mathbf{y},\mathbf{z}))\) 表示将 \(f_{i}\) 代入 \(X_{i}, g_{i}\)、将 \(Y_{i}, h_{i}\) 代入 \(Z_{i}\)、并按 \(1\leqslant i\leqslant n\) 的范围在 u 中得到的形式幂级数。对于每个系统 \(\alpha=(\alpha_{1},\ldots,\alpha_{n})\)（它由 n 个满足 \(\geqslant0\) 的整数构成），令 \(x^{\alpha}\) 表示单项式 \(X_{1}^{\alpha_{1}}\ldots X_{n}^{\alpha_{n}}\)；\(y^{\alpha}\) 和 \(z^{\alpha}\) 的定义类似。令 \(\varepsilon_{i}\) 表示这样的系统 \(\alpha\)，其满足 \(\alpha_{j}=0\)，条件为 \(j\neq i, \alpha_{i}=1\)。我们将系统 \((0,\ldots,0)\) 记为 e；它由 n 个 \(A_{x,y,z}\) 中的元素组成。

\begin{enumerate}
\def\labelenumi{(\alph{enumi})}
\item
  K 上的形式群律（或者滥用术语称 K 上的形式群）是维数为 n 的 \(G = \mathbf{f}(\mathbf{x}, \mathbf{y})\) 中 n 个元素的系统 \(A_{x,y}\)，满足以下性质：(1) \(\mathbf{f}(\mathbf{x}, \mathbf{f}(\mathbf{y}, \mathbf{z})) = \mathbf{f}(\mathbf{f}(\mathbf{x}, \mathbf{y}), \mathbf{z})\)；(2) \(\mathbf{f}(\mathbf{e}, \mathbf{y}) = \mathbf{y}\)、\(\mathbf{f}(\mathbf{x}, \mathbf{e}) = \mathbf{x}\)。证明必有 \(f_{i}(\mathbf{x}, \mathbf{y}) = \mathrm{X}_{i} + \mathrm{Y}_{i} + g_{i}(\mathbf{x}, \mathbf{y})\)，其中 \(g_{i}\) 只含次数 \(\geqslant 2\) 的单项式，且每个单项式都至少含有一个 \(X_{j}\) 和一个 \(Y_{j}\)。证明存在唯一一个由 \(\mathbf{h}(\mathbf{x})\) 中 n 个元素组成的系统 \(A_{x}\)，使得 \(\mathbf{f}(\mathbf{x}, \mathbf{h}(\mathbf{x})) = \mathbf{f}(\mathbf{h}(\mathbf{x}), \mathbf{x}) = \mathbf{e}\)（《代数学》，第 IV 章第 5 节第 9 号命题 10）。若 \(\mathbf{f}(\mathbf{y}, \mathbf{x}) = \mathbf{f}(\mathbf{x}, \mathbf{y})\)，则称 G 为交换的。
\item
对于所有 \(u \in \mathbf{A}_{\mathbf{x}}\)，令 \(L_{\mathbf{y}}u\) 表示 \(u(\mathbf{f}(\mathbf{y},\mathbf{x}))\) 中的元素 \(\mathbf{A}_{\mathbf{x},\mathbf{y}}\)。\(\mathbf{A}_{\mathbf{x}}\) 的导子 D 可规范地延拓为 \(\mathbf{A}_{\mathbf{x},\mathbf{y}}\) 的导子（同样记为 D，这里滥用记号），约定 \(\mathrm{D}(\mathrm{Y}_i) = 0\) 对所有 \(i\) 成立（《代数学》，第 IV 章第 5 节第 8 号命题 6）。若 \(L_{\mathbf{y}}D = DL_{\mathbf{y}}\)，则称 D 关于所讨论的形式群 G 左不变。记 \(\mathrm{D}^{(\mathrm{e})}\) 为从 \(\mathbf{A}_{\mathbf{x},\mathbf{y}}\) 到 \(\mathbf{A}_{\mathbf{y}}\) 的线性映射，其定义为 \(\mathrm{D}^{(\mathrm{e})}(u) = (\mathrm{D}u)(\mathbf{e},\mathbf{y})\)；它将 \(\mathbf{A}_{\mathbf{x}}\) 映入 K，并由其在 \(\mathbf{A}_{\mathbf{x}}\) 上的限制确定。证明 D 左不变的充要条件是对所有 \(v \in \mathbf{A}_{\mathbf{x}}\)，都有 \(\mathrm{D}^{(\mathrm{e})}(L_{\mathbf{y}}v) = (\mathrm{D}v)(\mathbf{y})\)。
\end{enumerate}

令 \(D_{i}\) 为 \(\partial/\partial X_{i}\) 上的导子 \(A_{x}\)（\(1 \leqslant i \leqslant n\)）；存在唯一一个左不变导子 \(T_{i}\) 满足 \(T_{i}^{(e)} = D_{i}^{(e)}\)。证明 \(T_{i}\) 在 K 上线性无关，并且每个左不变导子都是系数属于 K 的 \(T_{i}\) 的线性组合。对于括号 {[}D, D'{]} 和 p-映射 \(D \mapsto D^{p}\)（当 p \textgreater{} 0 时），推出左不变导子所成的集合 g 是一个李代数（当 p \textgreater{} 0 时是李 p-代数），称为形式李群 G 的李代数。

\begin{enumerate}
\def\labelenumi{(\alph{enumi})}
\setcounter{enumi}{2}
\tightlist
\item
  证明，对每个形式幂级数 \(u \in \mathbf{A}_{\mathbf{x}}\) :
\end{enumerate}

\[
u (\mathbf {f} (\mathbf {x}, \mathbf {y})) = u (0) + \sum_ {i = 1} ^ {n} Y _ {i} T _ {i} u + o _ {2} (\mathbf {x}, \mathbf {y})\tag{1}
\]

其中级数 \(o_{2}\)中的所有单项式次数均为 \(\geqslant2\)，关于 \(Y_{j}\)。推出

\[
\begin{array}{r l} {u (\mathbf {f} (\mathbf {x}, \mathbf {f} (\mathbf {y}, \mathbf {z}))) = u (0) + \sum_ {i = 1} ^ {n} (\mathrm{Y} _ {i} + \mathrm{Z} _ {i}) \mathrm{T} _ {i} u} & \\ {+ \sum_ {i, j} \mathrm{Y} _ {i} \mathrm{Z} _ {j} ((\mathrm{T} _ {i} \mathrm{T} _ {j}) (u)) + o _ {3} (\mathbf {x}, \mathbf {y}, \mathbf {z})} & \end{array}\tag{2}
\]

其中级数 \(o_{3}\)中的所有单项式次数均为 \(\geqslant3\)，关于集合 \(Y_{i}\) 和 \(Z_{i}\)。推出：

\[
u (\mathbf {f} (\mathbf {x}, \mathbf {y})) - u (\mathbf {f} (\mathbf {y}, \mathbf {x})) = \sum_ {i <   j} (X _ {i} Y _ {j} - X _ {j} Y _ {i}) ([ T _ {i}, T _ {j} ] ^ {(e)} (u)) + o _ {3} ^ {\prime} (\mathbf {x}, \mathbf {y})\tag{3}
\]

其中 \(o_3'\)中的所有项的总次数均为 \(\geqslant 3\)。证明，若 \(G\) 是交换的，\(g\) 是交换的。

\begin{enumerate}
\def\labelenumi{(\alph{enumi})}
\setcounter{enumi}{3}
\tightlist
\item
  设 G' 是另一个维数为 m 的形式群，其群律为 f'。 从 G 到 G' 的形式同态是一个系统 \(\mathbf{F} = (\mathrm{F}_{j}(\mathbf{x}))_{1 \leqslant j \leqslant m}\)，其元素属于 \(\mathbf{A}_{\mathbf{x}}\)，满足 \(\mathbf{f}'(\mathbf{F}(\mathbf{x}), \mathbf{F}(\mathbf{y})) = \mathbf{F}(\mathbf{f}(\mathbf{x}, \mathbf{y}))\) 。对于每个左不变导子 \(D \in g\) ，证明 \(v \mapsto D(v \circ F)\) 是属于李代数 \(g'\) 的左不变导子。 若记为 \(\mathbf{F}^*(D)\) , \(\mathbf{F}^*\)  是同态（当 p \textgreater{} 0 时为 p-同态），从 g 到 \(g'\) 。若 \((\mathrm{T}_{1})_{1 \leqslant i \leqslant n}, (\mathrm{T}'_{j})_{1 \leqslant j \leqslant m}\) 是 g 和 \(g'\) 的基，满足 \(T_{i}^{(e)} = D_{i}^{(e)}\)，且 \(T'_{j}{}^{(e)} = D'_{j}{}^{(e)}\) ，证明 \(F^*\) 关于这些基的矩阵为矩阵 \((\partial F_{j}/\partial X_{i})_{0}\)（\(\partial F_{j}/\partial X_{i}\)  的常数项；其中 i 是行的指标，j 是列的指标）。
\end{enumerate}

\begin{enumerate}
\def\labelenumi{\arabic{enumi}.}
\setcounter{enumi}{24}
\tightlist
\item
  在 \(n\) 个变量上，域 \(\mathbf{K}\)，记为 \(\mathbf{GL}(n, \mathbf{K})\)（无混淆时），是维数为 \(n^2\)，其群律由 \(f_{ij}(\mathbf{u}, \mathbf{v}) = -\delta_{ij} + \sum_{k=1}^{n} (\delta_{ik} + \mathrm{U}_{ik})(\delta_{kj} + \mathrm{V}_{kj})\)（\(\delta_{ij}\) 为克罗内克指标；一般定义中练习 24 的 \(3n^2\) 个不定元在这里记为 \(\mathrm{U}_{ij}, \mathrm{V}_{ij}, \mathrm{W}_{ij}\)，两个指标从 1 变化到 \(n\)）。证明，若 \(\mathrm{D}_{ij} = \partial/\partial \mathrm{U}_{ij}\)，满足该条件的左不变导子 \(\mathrm{X}_{ij}\) 满足 \(\mathrm{X}_{ij}^{(e)} = \mathrm{D}_{ij}^{(e)}\) 由下式给出
\end{enumerate}

\[
\mathrm{X} _ {i j} = (1 + \mathrm{U} _ {i i}) \mathrm{D} _ {i j} + \sum_ {k \neq i} \mathrm{U} _ {k i} \mathrm{D} _ {k j}.
\]

 的李代数 \(\mathbf{GL}(n,\mathbf{K})\) 被识别为 \(\mathfrak{gl}(n,\mathbf{K})\)，将 \(X_{ij}\) 与规范基中的元素 \(E_{ij}\)（使用练习 24 的公式 (3)）。若 K 的特征为 p \textgreater{}，\(X_{ii}^{p} = X_{ii}\)，且 \(X_{ij}^{p} = 0\) 当 \(i \neq j\)。

¶ 26. (a) 设 K 为特征为 \(\neq 2\)  的域，E 为 \(n\)  维 K 向量空间，\(\Phi\) 是 E 上的非退化对称双线性型。设 E 关于 \(\Phi\) 具有一组正交基，令 \(R\) 表示 \(\Phi\) 关于该基的每个矩阵 \(U\) 属于正交群 \(\mathbf{G}(\Phi)\) 且满足 \(\det(I + U) \neq 0\) 可写成

\[
U = (I - R ^ {- 1} S) ^ {- 1} (I + R ^ {- 1} S),
\]

  其中 \(S = R(U - I)(U + I)^{-1}\) 是交错矩阵，满足 \(\det(I - R^{-1}S) \neq 0\)；反之，对每个矩阵 \(S\) 满足这些条件的矩阵 \(U = (I - R^{-1}S)^{-1}(I + R^{-1}S)\) 是 \(\mathbf{G}(\Phi)\) 的矩阵，满足 \(\det(I + U) \neq 0\) 。令 \(S_{ij}, S_{ij}', S_{ij}''\) 为 \(3n(n - 1)/2\) 个不定元（ \(1 \leqslant i < j \leqslant n\) )

并令 L 为包含形式幂级数环 \(K[[S_{ij}, S'_{ij}, S''_{ij}]]\)。若以 \(S, S', S''\) 表示矩阵

\[
\sum_ {i <   j} \mathrm{S} _ {i j} (E _ {i j} - E _ {j i}), \qquad \sum_ {i <   j} \mathrm{S} _ {i j} ^ {\prime} (E _ {i j} - E _ {j i}), \qquad \sum_ {i <   j} \mathrm{S} _ {i j} ^ {\prime \prime} (E _ {i j} - E _ {j i}),
\]

则 \(\operatorname{det}(I - R^{-1}S) \neq 0\)，以及对 \(S'\) 和 \(S''\) 的类似结论也成立；若记

\[
U = (I - R ^ {- 1} S) ^ {- 1} (I + R ^ {- 1} S), \quad U ^ {\prime} = (I - R ^ {- 1} S ^ {\prime}) ^ {- 1} (I + R ^ {- 1} S ^ {\prime}),
\]

则还有 \(\operatorname{det}(I + UU') \neq 0\) . 记

\[
F (S, S ^ {\prime}) = R \left(U U ^ {\prime} - I\right) \left(U U ^ {\prime} + I\right) ^ {- 1} = \left(f _ {i j} \left(S, S ^ {\prime}\right)\right);
\]

其中 \(f_{ij}(S, S')\) 属于形式幂级数环 \(K[[S_{ij}, S_{ij}']]\)，并将 \(F(S, S')\) 看作系数属于 K 上 n 阶矩阵代数的形式幂级数，可写为：

\[
F (S, S ^ {\prime}) = S + S ^ {\prime} + o _ {2} (S, S ^ {\prime})
\]

其中，在矩阵\(o_2(S, S')\)的元素中，各项的次数关于\(\geqslant 1\)均为\(S_{ij}\)，且次数关于\(\geqslant 1\)均为\(S_{ij}'\)；同理：

\[
U = I + 2 R ^ {- 1} S + o _ {2} ^ {\prime} (S)
\]

 的元素是 \(o_2'(S)\) 为阶数 \(\geqslant 2\)。证明 \(F(S, S')\) 是 \(K\) 上的形式群律，维数为 \(n(n - 1)/2\) ；该形式群称为形式正交群（对应于 \(\Phi\) ），并记为 \(G(\Phi)\)（这里滥用记号）。

\begin{enumerate}
\def\labelenumi{(\alph{enumi})}
\setcounter{enumi}{1}
\tightlist
\item
  若记\(M(S) = (m_{ij}(S)) = (I - R^{-1}S)^{-1}(I + R^{-1}S)\) ，\(M\) 是从 \(\mathbf{G}(\Phi)\) 到 \(\mathbf{GL}(n, \mathbf{K})\)  的形式同态。令 \(\mathfrak{o}(\Phi)\) 为 \(\mathbf{G}(\Phi)\)，并令 \((\mathrm{T}_{ij})_{i < j}\) 表示该代数的一组基，满足 \(\mathrm{T}_{ij}^{(e)} = \mathrm{D}_{ij}^{(e)}\)（参见练习 24 (d)）；若元素 \(D = \sum_{i < j} c_{ij} T_{ij}\) 属于 \(\mathfrak{o}(\Phi)\) 被识别为矩阵 \(D = \sum_{i < j} c_{ij}(E_{ij} - E_{ji})\) ，证明（按练习 24 (d) 的记号，并使用练习 25 中将...的李代数与 \(\mathbf{GL}(n, \mathbf{K})\) 识别为 \(\mathfrak{gl}(n, \mathbf{K})\)）:
\end{enumerate}

\[
M ^ {*} (D) = 2 R ^ {- 1} D.
\]

  推出，\(M^*\) 是 \(\mathfrak{o}(\Phi)\) 到由 \(\mathfrak{gl}(n, \mathbf{K})\) 所成的子代数，矩阵 \(X\) 满足 \({}^{t}XR + RX = 0\)（对每个有序对 \((a, b)\) 个元素 \(E\) ，将其识别为单列矩阵，考察 \(\Phi(a + M(S).a, b + M(S).b) = {}^{t}a.(I + {}^{t}M(S))R(I + M(S)).b\) 作为取常数项后的形式幂级数，并利用由 \(X \in \mathfrak{gl}(n, \mathbf{K})\) 满足 \({}^{t}XR + RX = 0\) 的子代数维数为 \(n(n - 1)/2\)）。

\begin{enumerate}
\def\labelenumi{(\alph{enumi})}
\setcounter{enumi}{2}
\tightlist
\item
  类似地定义任意
\end{enumerate}

（交换）域上，并证明其李代数被识别为由 \(\mathfrak{gl}(2n,\mathbf{K})\) 所成的子代数，其中矩阵 X 满足

\[
{ } ^ { t } X A + A X = 0 , \quad \text {   where   } \quad A = \left( \begin{array} { c c } 0 & I _ { n } \\ - I _ { n } & 0 \end{array} \right) .
\]

\begin{enumerate}
\def\labelenumi{\arabic{enumi}.}
\setcounter{enumi}{26}
\tightlist
\item
  设 K 为特征为 p \textgreater{} 为 0，G 为由 \(f_{1}(\mathbf{x}, \mathbf{y}) = \mathrm{X}_{1} + \mathrm{Y}_{1} + \mathrm{X}_{1}\mathrm{Y}_{1}, f_{2}(\mathbf{x}, \mathbf{y}) = \mathrm{X}_{2} + \mathrm{Y}_{2}(1 + \mathrm{X}_{1}^{p})\)。证明 G 不交换，但其李代数是交换的。
\end{enumerate}

\exercisesfor{2}

\begin{enumerate}
\def\labelenumi{\arabic{enumi}.}
\item
  采用定义 1 和第 6 号中的记号 T、J、\(T_{n}\)。令 t 为 T 的 p 阶齐次张量，且 \(\sigma\) 是 \(\{1, 2, \ldots, p\}\) 的一个排列。则 \(t - \sigma t \in J + T_{p-1}\)。（将其归结为 \(\sigma\) 是两个相邻整数的换位。）
\item
  设 \(\mathbf{K}\) 是一个域。令 \(\mathfrak{g}\) 为李 \(\mathbf{K}\) 代数，其 \(\mathbf{U}\) 包络代数。
\end{enumerate}

\begin{enumerate}
\def\labelenumi{(\alph{enumi})}
\item
  设 \(u, v\) 属于 U。若 \(u\) 的滤过是 \(> n\)，且 \(v\) 的滤过是 \(> p\)，则 \(uv\) 的滤过为 \(> n + p\)（使用定理 1）。
\item
  由 (a) 推出，U 的可逆元素只有标量。
\item
  由 (b) 推出，U 的 Jacobson 根为零。
\end{enumerate}

\begin{enumerate}
\def\labelenumi{\arabic{enumi}.}
\setcounter{enumi}{2}
\item
  设 K 是一个域。令 g 为李 K-代数，U 为其包络代数。 U 的左正则表示对应于 g 上的表示 \(\rho\)，作用在 U 空间上。证明集合 \(U_{+}\) 是 U 中不含常数项的元素所成的集合，在 \(\rho\) 下稳定，且 \(U_{+}\) 在 U 中没有关于 \(\rho\) 的补空间。特别地，\(\rho\) 不是半单的。（将在 § 6 的定理 2 中再次讨论。）
\item
  设 K 是一个域。
\end{enumerate}

\begin{enumerate}
\def\labelenumi{(\alph{enumi})}
\item
  验证存在一个具有基 \((x, y, z)\)，满足 \([x, y] = z\) , \([x, z] = [y, z] = 0\)。令 U 为 g 的包络代数。证明 U 的中心是由 1 和 z 生成的子代数。（对每个 \(\xi \in g\)，考虑将 ad 扩张为 g 的对称代数 S 上的导子，\(\xi\)，寻找被这些导子消去的 S 中元素；然后应用第 8 号末尾的注记。）
\item
  验证存在一个具有基 \((x, y, z)\)，满足 \([x, y] = y\) , \([x, z] = z\) , \([y, z] = 0\)。证明 g 的包络代数的中心退化为标量。（方法相同。）
\end{enumerate}

\begin{enumerate}
\def\labelenumi{\arabic{enumi}.}
\setcounter{enumi}{4}
\item
  设 \(\mathbf{K}\) 是特征为 \(p > 0\)（\(p\) 为素数）。令 \(\mathfrak{g}\) 为 \(\mathbf{K}\)，具有一组基，并规范地识别为其包络代数 U 的子模。对于 \(\phi\) 的 K-模 \(\mathfrak{g}\) 成为 \(p\)-映射的充要条件是 \(x \mapsto x^p - \phi(x)\) 是关于 \(\lambda \mapsto \lambda^p)\) 的半线性映射，它将 \(\mathfrak{g}\) 映入 U 的中心。推出，若 \((b_{\lambda})\) 是 \(\mathfrak{g}\)，则存在 \(p\) -映射定义在 \(\mathfrak{g}\) 上的充要条件是，对所有 \(\lambda\)，存在 \(c_{\lambda} \in \mathfrak{g}\) 满足 \((\operatorname{ad} b_{\lambda})^p = \operatorname{ad} c_{\lambda}\)；若如此，则存在唯一一个 \(p\) -映射 \(x \mapsto x^{[p]}\) 满足 \(b_{\lambda}^{[p]} = c_{\lambda}\) 对所有 \(\lambda\)。
\item
  设 g 是环 K 上的李 p-代数，满足 \(pK = \{0\}\)（p 为素数），U 为其包络代数，\(\sigma\) 是将 g 映入 U 的典范映射。令 J 为 U 的双边理想，由下列元素生成：\((\sigma(x))^{p} - \sigma(x^{[p]})\)，其中 x 遍历 g。结合代数 \(\tilde{U} = U/J\) 称为 g 的限制包络代数。映射 \(\sigma\) 在取商后定义映射 \(\tilde{\sigma}\)（称为典范映射）将 g 映入 \(\tilde{U}\) ，它是 p-同态（当 \(\tilde{U}\) 被视为李 p-代数时）。
\end{enumerate}

\begin{enumerate}
\def\labelenumi{(\alph{enumi})}
\item
  代数 \(\tilde{\mathbf{U}}\) 与映射 \(\tilde{\sigma}\) 是以下泛映射问题的解：对每个 \(p\) -同态 \(f\)，将 \(\mathfrak{g}\) 映入结合代数 \(\mathbf{B}\) 上的 \(\mathbf{K}\)（视为李 \(p\)-代数），存在唯一一个 K-同态 \(f'\)，将 \(\tilde{\mathbf{U}}\) 映入 \(\mathbf{B}\)（连同它们的结合代数结构）满足 \(f = f' \circ \tilde{\sigma}\)。
\item
  证明，若 \((b_{\lambda})_{\lambda \in \Lambda}\) 是 \(\mathfrak{g}\)（其中 \(\Lambda\) 完全有序），\(\tilde{\sigma}\) 是单射，且若 \(x\) 与 \(\tilde{\sigma}(x)\) 在 \(x\in \mathfrak{g}\) 中视为同一对象，则元素 \(\Pi_{\lambda}b_{\lambda}^{\nu_{\lambda}}\)（其中 \(\lambda\) 按递增顺序排列，\(\nu_{\lambda}\) 除有限多个外均为零，且 \(0\leqslant \nu_{\lambda} < p\) 对所有 \(\lambda\)）构成 \(\tilde{\mathbf{U}}\) 的一组基。（规范地将 \(\mathfrak{g}\) 与 \(\mathrm{U}\) 识别为子模，记 \(\phi (x) = x^{p} - x^{[p]}\) 对所有 \(x\in g\)。（对每个复合指标 \(\alpha = (\alpha_{\lambda})\in \mathbf{N}^{(\Lambda)}\)，令 \(\alpha_{\lambda} = \beta_{\lambda} + p\gamma_{\lambda}\) 满足 \(0\leqslant \beta_{\lambda} < p\)，并令
\end{enumerate}

\[
\mathrm{T} _ {\alpha} = (\Pi_ {\lambda} b _ {\lambda} ^ {\beta_ {\lambda}}) (\Pi_ {\lambda} (\phi (b _ {\lambda})) ^ {\gamma_ {\lambda}}).
\]

证明 \(T_{\alpha}\) 构成 U 的一组基；而满足 \(T_{\alpha}\) 的 \(\gamma = (\gamma_{\lambda}) \neq 0\) 构成 J 的一组基；注意 \(\phi(b_{\lambda})\) 属于 U 的中心。）

\begin{enumerate}
\def\labelenumi{\arabic{enumi}.}
\setcounter{enumi}{6}
\tightlist
\item
  设 g 是环 K 上的李 p-代数，满足 \(pK = \{0\}\)（p 为素数）。g 的导子 D 称为 p-导子，如果 \(\mathrm{D}(x^{p}) = (\mathrm{ad} x)^{p-1}\)。对所有 \(x \in g\)。每个内导子都是 p-导子。
\end{enumerate}

\begin{enumerate}
\def\labelenumi{(\alph{enumi})}
\item
  若 \(\mathbf{L}\) 是结合 K-代数，则 \(\mathbf{L}\) 的每个导子都是 \(p\)-导子，当 \(\mathbf{L}\) 被视为李 \(p\)-代数时。（使用第 1 节练习 19 的公式 (2)。）
\item
  设 g 有一组基。g 的导子为 p-导子的充要条件是它可延拓为 g 的限制包络代数的导子。推出，g 的 p-导子构成 g 的导子所成 p-代数的李 p-子代数。
\item
  若 D 是 \(p\) -导子，作用于 \(\mathfrak{g}\) ，则 D 的核是 \(p\) -子代数。\(\mathfrak{g}\) .
\item
  对于 g 的每个导子 D，\(\mathrm{D}(x^{p}) - (\mathrm{ad} x)^{p-1}\) . \(\bar{D}x\) 属于 g 的中心，对所有 \(x \in g\)（使用第 1 节练习 19 的公式 (2)，应用于 \(\mathrm{L} = \mathcal{L}(\mathfrak{g})\)）。
\end{enumerate}

\begin{enumerate}
\def\labelenumi{\arabic{enumi}.}
\setcounter{enumi}{7}
\item
  证明，若模 \(\mathfrak{g}\) 是单生成模的直和，定理 1 仍然成立。（在证明中用 g 的对称代数替换模 P。）
\item
  令\(k\)为含两个元素的域，\(\mathbf{V}\)为向量空间\(k^3\)，\((x_1, x_2, x_3)\)为其规范基，且\(\mathbf{K}\)为\(\mathbf{V}\)的外代数；它是一个定义在\(k\)上的 8 维交换代数。令\(\mathfrak{g}\)为具有基\((e_1, e_2, e_3, e_{12}, e_{13}, e_{23})\)的李 K-代数，使得
\end{enumerate}

\[
\left[ e _ {1}, e _ {2} \right] = \left[ e _ {2}, e _ {1} \right] = e _ {1 2} \quad \left[ e _ {1}, e _ {3} \right] = \left[ e _ {3}, e _ {1} \right] = e _ {1 3}, \quad \left[ e _ {2}, e _ {3} \right] = \left[ e _ {3}, e _ {2} \right] = e _ {2 3}
\]

且其他括号均为零。令 \(\mathfrak{h}\) 为 \(\mathfrak{g}\) 的理想，由 \(u = x_{1}e_{1} + x_{2}e_{2} + x_{3}e_{3}\) 生成。作为模，\(\mathfrak{h}\) 由 \(u\)、\([u, e_1] = x_2e_{12} + x_3e_{13}\) , \([u, e_2] = x_1e_{12} + x_3e_{23}\) , \([u, e_3] = x_1e_{13} + x_2e_{23}\)。令

\[
v = x _ {1} x _ {2} e _ {1 2} + x _ {1} x _ {3} e _ {1 3} + x _ {2} x _ {3} e _ {2 3}.
\]

\[
\phi \left(e _ {1}\right) = \phi \left(e _ {2}\right) = \phi \left(e _ {3}\right) = 0, \phi \left(e _ {1 2}\right) = x _ {3}, \phi \left(e _ {1 3}\right) = x _ {2}, \phi \left(e _ {2 3}\right) = x _ {1}.
\]

\begin{enumerate}
\def\labelenumi{(\alph{enumi})}
\setcounter{enumi}{1}
\item
  设 \(f\) 为从 \(\mathfrak{g}\) 到结合 K-代数的线性映射，满足 \(f([x,y]) = f(x)f(y) - f(y)f(x)\) 对所有 \(x, y\) 中的 \(\mathfrak{g}\)。证明 \(f(v) = f(u)^2\)。
\item
  由 (a) 和 (b) 推出，李代数 \(\mathfrak{g} / \mathfrak{h}\) 到其包络代数的典范映射不是单射。
\end{enumerate}

\begin{enumerate}
\def\labelenumi{\arabic{enumi}.}
\setcounter{enumi}{9}
\tightlist
\item
  设 \(\mathfrak{g}\) 是域上的有限维李代数，U 是其包络代数，\(\mathbf{U}_n\) 是 U 中滤过 \(\leqslant n\) 的元素所成的集合。
\end{enumerate}

\begin{enumerate}
\def\labelenumi{(\alph{enumi})}
\item
  设 \(x \in \mathrm{U}_n, y \in \mathrm{U}_n\) 是两个非零元素。证明存在 \(u \in \mathrm{U}, v \in \mathrm{U}\) 使得 \(ux = vy\)。（比较 \(\dim(\mathrm{U}_p x) = \dim \mathrm{U}_p\)，\(\dim(\mathrm{U}_p y) = \dim \mathrm{U}_p\) 和 \(\dim \mathrm{U}_{p+n}\)。推出 \(\mathrm{U}_p x \cap \mathrm{U}_p y \neq \{0\}\) 对于 \(p\) 足够大时成立。）
\item
  证明 U 存在一个左商域（《代数学》，第 I 章第 9 节练习 8），同时也是右商域。
\end{enumerate}

\exercisesfor{3}

\begin{enumerate}
\def\labelenumi{\arabic{enumi}.}
\item
  设 g 为李代数，\(\rho\) 是 g 在 K-模 M 上的表示，\(\sigma\) 是 g 在 M 的张量代数上的相应表示。证明，对称张量子模和斜对称张量子模关于 \(\sigma\) 稳定。
\item
  设 g 为李代数，\(\rho\) 是 g 在 K-模 M 上的表示，\(\sigma\) 是 g 在 Q = L(M, M; M) 上的相应表示。对于 \(f \in Q\) 关于 \(\sigma\)  的不变性充要条件是 \(\rho(x)\) 是关于 f 所定义乘法的 M 的导子。
\item
  设 V 是完美域 K 上的有限维向量空间。
\end{enumerate}

\begin{enumerate}
\def\labelenumi{(\alph{enumi})}
\tightlist
\item
   的恒等表示 \(\mathfrak{gl}(\mathrm{V})\) 在 \(\mathfrak{gl}(\mathrm{V})\) 上规范地定义表示，作用于 \(\bigotimes_{q}^{p}\mathrm{V}\) 和 \(\bigotimes_{q}^{q}\mathrm{V}^{*}\)，从而表示 \(u\mapsto u_q^p\) \(\mathfrak{gl}(\mathrm{V})\) 作用于
\end{enumerate}

\(V_{q}^{p} = \left( \bigotimes_{i=1}^{p} V \right) \otimes \left( \bigotimes_{i=1}^{q} V^{*} \right)\)。证明 \(u_{q}^{p}\) 若 \(u\) 是半单的。（延拓基域，将其归结为 \(K\) 为代数闭域的情形。）证明 \(u_{q}^{p}\) 若 \(u\) 是幂零的。推出，若 \(s\) 与 \(n\) 是 \(u, u_{q}^{p}\) 和 \(n_{q}^{p}\) 的半单与幂零分量，则 \(u_{q}^{p}\) 的半单与幂零分量。

\begin{enumerate}
\def\labelenumi{(\alph{enumi})}
\setcounter{enumi}{1}
\item
  由 (a) 推出，若 \(V_q^p\) 中的一个元素被 \(u_q^p\)  消去，则它也被 \(s_q^p\) 和 \(n_q^p\) 消去。
\item
  由 (b) 和练习 2 推出，若 V 具有不必结合的代数结构且 u 是 V 的导子，则 s 和 n 都是 V 的导子。
\end{enumerate}

\begin{enumerate}
\def\labelenumi{\arabic{enumi}.}
\setcounter{enumi}{3}
\tightlist
\item
  设 K 是一个域。令 g 为李 K-代数。
\end{enumerate}

\begin{enumerate}
\def\labelenumi{(\alph{enumi})}
\item
  设 M、N 为 g-模，其中 N 是 K 上的有限维模。令 \((f_{j})_{1 \leqslant j \leqslant n}\) 为 N 在 K 上的一组基。若存在 M 中的元素 \(e_{1}, \ldots, e_{n}\)，不全为零，使得 \(\sum_{i=1}^{n} e_i \otimes f_i\) 在 \(\mathbf{M} \otimes \mathbf{N}\)  中不变，则由 \(\mathbf{M}_1\) 生成的 \(\mathbf{M}\) 是 \(e_i\)，关于 \(x_{\mathbf{M}}\) 稳定；若进一步 \(\mathbf{N}\) 是单的，则 g 在 \(\mathfrak{g}\) 上的表示 \(\mathbf{M}_1\) 对偶于 g 在 \(\mathfrak{g}\) 上的表示。\(\mathbf{N}\) .
\item
  设 \(M_{1}\) 、\(M_{2}\)、N、\(N^{*}\) 是 K 上的有限维 g-模。设 \(M_{1}\) 与 \(M_{2}\) 是单的，且 g 在 N 和 \(N^{*}\) 上的表示互为对偶。对于 g-模 \(M_{1}\) 要同构于 \(M_{2} \otimes N\)  的一个 g-子模的充要条件是 \(M_{2}\) 同构于 \(M_{1} \otimes N^{*}\)。（使用命题 4；考察 g 在 \(M_{1}^{*} \otimes M_{2} \otimes N\) 上的表示，并将 (a) 应用于后者的对偶表示。）
\end{enumerate}

\begin{enumerate}
\def\labelenumi{\arabic{enumi}.}
\setcounter{enumi}{4}
\tightlist
\item
  设 K 是一个特征为 0 的域。令 V 为有限维 K-向量空间，\(\mathfrak{g} = \mathfrak{sl}(V)\)，U 为 \(\mathfrak{g}\) 的包络代数，\(U_n\) 是 U 中滤子不超过 \(\leqslant n\) 的元素集合；\(U^n\) 是 U 中阶数为 \(n\) 的、在 \(\mathfrak{g}\) 上定义的齐次对称张量的像，且\((\alpha_1, \ldots, \alpha_m)\) 是向量空间 g 的一组基。令 \(W_{s} = \bigotimes^{s} V\)。g 的恒等表示在 \(W_{s}\) 上定义了 g 的一个表示，并将其延拓为同态 \(\pi_s\)，从 U 到代数 \(M_s = \mathcal{L}(W_s)\)。
\end{enumerate}

\begin{enumerate}
\def\labelenumi{(\alph{enumi})}
\tightlist
\item
  令\(M_{s}^{\prime}\)为由\(M_{s}\)中形如\(\beta_{1} \otimes \beta_{2} \otimes \cdots \otimes \beta_{s}\)的元素生成的子空间，其中\(\beta_{i} \in \mathcal{L}(V)\)，且\(\beta_{i} = 1\)对至少一个 i 等于 1。令 z 为\(U_{s}\)中的一个元素，并记其\(z_{0} = \sum_{1 \leqslant i_{1}, \ldots, i_{s} \leqslant m} a_{i_{1} \ldots i_{s}} \alpha_{i_{1}} \ldots \alpha_{i_{s}}\)为\(U^{s}\)中的分量；其中\(a_{i_{1} \ldots i_{s}}\)关于指标置换是对称的。证明
\end{enumerate}

\[
\pi_ {s} (z) \equiv s! \sum_ {1 \leqslant i _ {1}, \dots , i _ {s} \leqslant m} a _ {i _ {1} \dots i _ {s}} \alpha_ {i _ {1}} \otimes \dots \otimes \alpha_ {i _ {s}} \pmod {\mathrm{M} _ {s} ^ {\prime}}.
\]

\begin{enumerate}
\def\labelenumi{(\alph{enumi})}
\setcounter{enumi}{1}
\tightlist
\item
  由 (a) 推出，对任意\(z \in \mathbf{U}\)，存在一个有限维表示\(\pi\)，其作用于\(\mathbf{U}\)满足\(\pi(z) \neq 0\)不为 0。（关于本题的后续内容，参见 § 7，习题 3。）
\end{enumerate}

\begin{enumerate}
\def\labelenumi{\arabic{enumi}.}
\setcounter{enumi}{5}
\tightlist
\item
 设\(\mathbf{K}\)是一个域。令\(\mathfrak{g}\)为李 K-代数，\(x \mapsto x_{\mathbf{M}}\)是\(\mathfrak{g}\)的有限维表示，\((e_1, \ldots, e_n)\)是\(\mathbf{M}\)的一组基，而\((e_1^*, \ldots, e_n^*)\)是对偶基。
\end{enumerate}

\begin{enumerate}
\def\labelenumi{(\alph{enumi})}
\tightlist
\item
  若\(f\)是定义在\(\bigotimes^{r} \mathbf{M}\)上的线性形式，则：
\end{enumerate}

\[
f = \sum_ {1 \leqslant i _ {1} \dots , i _ {r} \leqslant n} f (e _ {i _ {1}} \otimes \dots \otimes e _ {i _ {r}}) (e _ {i _ {1}} ^ {*} \otimes \dots \otimes e _ {i _ {r}} ^ {*}).
\]

\begin{enumerate}
\def\labelenumi{(\alph{enumi})}
\setcounter{enumi}{1}
\tightlist
\item
  令\(\beta\)为定义在\(\mathbf{M}\)上且在\(\mathfrak{g}\)下不变的非退化双线性形式。令\((e_1', \ldots, e_n')\)为\(\mathbf{M}\)的一组基，满足\(\beta(e_i, e_j') = \delta_{ij}\)。令\(f\)为定义在\(\bigotimes^r \mathbf{M}\)上的不变线性形式。由 (a) 推出：
\end{enumerate}

\[
\sum_ {1 \leqslant i _ {1}, \dots , i _ {r} \leqslant n} f (e _ {i _ {1}} \otimes \dots \otimes e _ {i _ {r}}) (e _ {i _ {1}} ^ {\prime} \otimes \dots \otimes e _ {i _ {r}} ^ {\prime})
\]

是\(\otimes\) M 中的一个不变元，且与基\((e_i)\)的选择无关。

\begin{enumerate}
\def\labelenumi{(\alph{enumi})}
\setcounter{enumi}{2}
\item
  令 U 为 g 的包络代数，U* 为 U 的对偶空间。g 的伴随表示可延拓为表示\(x \mapsto x_{\mathrm{U}}\)：它在 U 上由\(x_{\mathrm{U}}u = xu - ux\)定义，对所有\(u \in U\)均成立。因此 U* 具有 g-模结构。对于 U* 中的元素 f，其为不变元的充要条件是\(f(uv) = f(vu)\)，对 U 中所有 u、v 均成立。
\item
  令\(\mathfrak{h}\)为\(\mathfrak{g}\)的有限维理想，\(\gamma\)为定义在\(\mathfrak{g}\)上的不变双线性形式，且其在\(\mathfrak{h}\)上的限制非退化。令\((y_i)_{1\leq i\leq n}, (y_j')_{1\leq j\leq n}\)为\(\mathfrak{h}\)的两组基，满足\(\gamma(y_i,y_j') = \delta_{ij}\)。令\(r\)为一个\(\geqslant 1\)的整数。令\(f\)为 U 上的线性形式，满足\(f(w) = f(vu)\)，对 U 中所有\(u,v\)均成立。由 (b) 和 (c) 推出：
\end{enumerate}

\[
\sum_ {1 \leqslant t _ {1}, \dots t _ {r} \leqslant n} f (y _ {t _ {1}} y _ {t _ {2}} \cdot \cdot \cdot y _ {t _ {r}}) y _ {t _ {1}} ^ {\prime} y _ {t _ {2}} ^ {\prime} \cdot \cdot \cdot y _ {t _ {r}} ^ {\prime}
\]

是 U 中心里与基\((y_{i})\)的选择无关的元素。特别地，这给出了 Casimir 元。

\begin{enumerate}
\def\labelenumi{(\alph{enumi})}
\setcounter{enumi}{4}
\tightlist
\item
  令\(\theta\)为\(\{1, \ldots, r\}\)的一个置换。用类似的论证证明：
\end{enumerate}

\[
\sum_ {1 \leqslant i _ {1}, \dots i _ {r} \leqslant n} f (y _ {i _ {\theta (1)}} y _ {i _ {\theta (2)}} \dots y _ {i _ {\theta (r)}}) y _ {i _ {1}} ^ {\prime} y _ {i _ {2}} ^ {\prime} \dots y _ {i _ {r}} ^ {\prime}
\]

是 U 中心里与基\((y_{i})\)的选择无关的元素。

\begin{enumerate}
\def\labelenumi{\arabic{enumi}.}
\setcounter{enumi}{6}
\item
  令 g 为复李代数，\(\beta\)为其 Killing 形式，\(g_{0}\)为将基域限制为 R 后由 g 得到的李代数。证明\(g_{0}\)的 Killing 形式是\(\beta\)的实部的两倍。
\item
  令\(\mathfrak{g}\)为实李代数。多线性形式
\end{enumerate}

\[
(x _ {1}, \dots , x _ {n}) \mapsto \operatorname{Tr} (\operatorname{ad} x _ {1} \circ \operatorname{ad} x _ {2} \circ \dots \circ \operatorname{ad} x _ {n})
\]

在\(\mathfrak{g}^n\)上是不变的。

\begin{enumerate}
\def\labelenumi{\arabic{enumi}.}
\setcounter{enumi}{8}
\item
  令 g 为李 K-代数，\(\rho\)和\(\rho'\)是 g 在 K-模 V、\(V'\)上的半单表示，且\(\phi\)是从 g-模 V 到 g-模\(V'\)的满同态。证明，\(\phi\)将 V 的不变元集合映到\(V'\)的不变元集合（使用命题 6）。
\item
  设 K 是一个域。令 g 为李 K-代数，\(M_{1}\)、\(M_{2}\)是两个不同构的有限维简单 g-模。若\(K_{1}\)是 K 的可分扩张，证明不存在简单的\(g_{(K_{1})}\)-模，它同时同构于\(g_{(K_{1})}\)-模\(M_{1(K_{1})}\)的一个子模，也同构于\(g_{(K_{1})}\)-模\(M_{2(K_{1})}\)的一个子模。（使用命题 4，并注意：\(\neq0\)在\(M_{1(K_{1})}^{*}\otimes_{K}M_{2(K_{1})}\)中表示一个不变元，意味着\(\neq0\)在\(M_{1}^{*}\otimes_{K}M_{2}\)中表示一个不变元（《代数》第二章 § 5，第 3 号，定理 1）。）
\item
  令\(\mathfrak{w}\)为 § 1 习题 21 中定义的李 \(p\)-代数。证明\(\mathfrak{w}\)的 Killing 形式为零。
\end{enumerate}

¶ 12. 设 K 是一个域。令 g 为李 K-代数，M 为 g-模。令\(C^{p}(\mathfrak{g}, \mathrm{M})\)表示从\(g^{p}\)到 M 的交错线性映射空间。记\(C^{0}(\mathfrak{g}, \mathrm{M}) = \mathrm{M}\)，并且当 p\textless{} 0 时，\(C^{p}(\mathfrak{g}, \mathrm{M}) = \{0\}\)。令\(C^{*}(\mathfrak{g}, \mathrm{M})\)为各个\(C^{p}(\mathfrak{g}, \mathrm{M})\)的直和。\(C^{*}(\mathfrak{g}, \mathrm{M})\)中的元素称为 g 取值于 M 的上链；\(C^{p}(\mathfrak{g}, \mathrm{M})\)中的元素称为 p 次上链。对于所有\(y \in g\)，记\(i(y)\)为\(C^{*}(\mathfrak{g}, \mathrm{M})\)的自同态；它将每个子空间\(C^{p}(\mathfrak{g}, \mathrm{M})\)映到\(C^{p-1}(\mathfrak{g}, \mathrm{M})\)，并且当 p\textgreater{} 0 时由下式给出：

\[
(i (y) f) \left(x _ {1}, \dots , x _ {p - 1}\right) = f \left(y, x _ {1}, \dots , x _ {p - 1}\right).\tag{1}
\]

有\(i(y)^{2}=0\)。

\begin{enumerate}
\def\labelenumi{(\alph{enumi})}
\tightlist
\item
  g 的伴随表示与 g 在 M 上的表示在从\(g^{p}\)到 M 的多线性映射空间上定义了 g 的一个表示。证明\(C^{p}(g, M)\)在此表示下稳定。令\(\theta\)为 g 在\(C^{*}(g, M)\)上的上述表示。证明：
\end{enumerate}

\[
\theta (x) i (y) - i (y) \theta (x) = i ([ x, y ])
\]

对 g 中所有 x、y。

\begin{enumerate}
\def\labelenumi{(\alph{enumi})}
\setcounter{enumi}{1}
\tightlist
\item
  证明存在唯一的自同态\(d\)，它作用于\(\mathrm{C}^*(\mathfrak{g},\mathrm{M})\)，将\(\mathrm{C}^p (\mathfrak{g},\mathrm{M})\)映到\(\mathrm{C}^{p + 1}(\mathfrak{g},\mathrm{M})\)，使得
\end{enumerate}

\[
d i (y) + i (y) d = \theta (y)\tag{3}
\]

对所有\(y \in \mathfrak{g}\)。（按上链次数作归纳。）证明，对\(f \in \mathrm{C}^p(\mathfrak{g}, \mathbf{M})\)：

\[
\begin{array}{r l} d f (x _ {1}, x _ {2}, \dots , x _ {p + 1}) & = \sum_ {i <   j} (- 1) ^ {i + j} f ([ x _ {i}, x _ {j} ], x _ {1}, \dots , \hat {x} _ {i}, \dots , \hat {x} _ {j}, \dots , x _ {p + 1}) \\ & \quad + \sum_ {i} (- 1) ^ {i + 1} (x _ {i}) _ {\mathrm{M}} f (x _ {1}, \dots , \hat {x} _ {i}, \dots , x _ {p + 1}) \end{array}
\]

（其中字母上方的符号\^{}表示省略该字母）。

\begin{enumerate}
\def\labelenumi{(\alph{enumi})}
\setcounter{enumi}{2}
\tightlist
\item
  证明，对所有\(y \in g\)，
\end{enumerate}

\[
d \theta (y) = \theta (y) d.\tag{4}
\]

（先利用 (2) 和 (3) 证明\(d\theta(y) - \theta(y)d\)与每个\(i(x)\)交换。然后按上链次数作归纳。）

\begin{enumerate}
\def\labelenumi{(\alph{enumi})}
\setcounter{enumi}{3}
\item
  证明\(d^2 = 0\)。（先利用 (3) 和 (4) 证明\(d^2\)与每个\(i(x)\)交换。然后按上链次数作归纳。）
\item
  \(d\)在\(\mathrm{C}^p (\mathfrak{g},\mathrm{M})\)上的限制的核为\(\mathrm{Z}^p (\mathfrak{g},\mathrm{M})\)，其中元素称为取值于 M 的\(p\)次上循环。\(d\)在\(\mathrm{C}^{p - 1}(\mathfrak{g},\mathrm{M})\)上的限制的像为\(\mathrm{B}^p (\mathfrak{g},\mathrm{M})\)，其中元素称为取值于 M 的\(p\)次上余边界。因此\(\mathrm{B}^p (\mathfrak{g},\mathrm{M})\subset \mathrm{Z}^p (\mathfrak{g},\mathrm{M})\)。商空间
\end{enumerate}

\[
\mathrm{Z} ^ {p} (\mathfrak {g}, \mathrm{M}) / \mathrm{B} ^ {p} (\mathfrak {g}, \mathrm{M}) = \mathrm{H} ^ {p} (\mathfrak {g}, \mathrm{M})
\]

称为 g 的 p 次、取值于 M 的上同调空间。\(\mathrm{H}^{p}(\mathfrak{g},\mathrm{M})\)的直和记为\(\mathrm{H}^{*}(\mathfrak{g},\mathrm{M})\)。证明\(\mathrm{H}^{0}(\mathfrak{g},\mathrm{M})\)可与 M 的不变元集合等同。令\(\phi\)为从 g-模 M 到 g-模 N 的同态。对所有\(f \in \mathrm{C}^{p}(\mathfrak{g},\mathrm{M})\)，\(\phi \circ f \in \mathrm{C}^{p}(\mathfrak{g},\mathrm{N})\)，因此\(\phi\)可延拓为 K-线性映射\(\phi'\)：\(\mathrm{C}^{*}(\mathfrak{g},\mathrm{M}) \to \mathrm{C}^{*}(\mathfrak{g},\mathrm{N})\)。证明\(\phi' \circ d = d \circ \phi'\)。由此推出\(\phi'\)定义了同态

\[
\tilde {\phi} \colon \mathrm{H} ^ {*} (\mathfrak {g}, \mathrm{M}) \to \mathrm{H} ^ {*} (\mathfrak {g}, \mathrm{N})
\]

其称为与\(\phi\)相关的同态。

\begin{enumerate}
\def\labelenumi{(\alph{enumi})}
\setcounter{enumi}{5}
\tightlist
\item
  令 L 为 M 的子 g-模，N 为商 g-模 M/L。规范同态\(L \stackrel{\iota}{\rightarrow} M \stackrel{p}{\rightarrow} N\)定义了同态：
\end{enumerate}

\[
\begin{array}{l} \mathrm{C} ^ {*} (\mathfrak {g}, \mathrm{L}) \xrightarrow {\iota^ {\prime}} \mathrm{C} ^ {*} (\mathfrak {g}, \mathrm{M}) \xrightarrow {p ^ {\prime}} \mathrm{C} ^ {*} (\mathfrak {g}, \mathrm{N}), \\ \mathrm{H} ^ {n} (\mathfrak {g}, \mathrm{L}) \xrightarrow {\iota^ {n}} \mathrm{H} ^ {n} (\mathfrak {g}, \mathrm{M}) \xrightarrow {p ^ {n}} \mathrm{H} ^ {n} (\mathfrak {g}, \mathrm{N}). \end{array}
\]

证明\(\iota'\)是单射，且其像为\(p'\)的核，并且\(p'\)是满射。对任意\(c \in \mathrm{H}^n(\mathfrak{g},\mathrm{N})\)，令\(z \in \mathrm{Z}^n(\mathfrak{g},\mathrm{N})\)为\(c\)的一个代表元，并令\(a \in \mathrm{C}^n(\mathfrak{g},\mathrm{M})\)满足\(p'(a) = z\)；证明\(da \in \mathrm{Z}^{n+1}(\mathfrak{g},\mathrm{L})\)，且其在\(\mathrm{H}^{n+1}(\mathfrak{g},\mathrm{L})\)中的类只依赖于\(c\)；将此类记为\(\delta^n c\)。证明以下同态序列：

\[
\begin{array}{r l} 0 \longrightarrow & \mathrm{H} ^ {0} (\mathfrak {g}, \mathrm{L}) \stackrel {{\iota ^ {0}}} {{\longrightarrow}} \mathrm{H} ^ {0} (\mathfrak {g}, \mathrm{M}) \stackrel {{p ^ {0}}} {{\longrightarrow}} \mathrm{H} ^ {0} (\mathfrak {g}, \mathrm{N}) \stackrel {{\delta^ {0}}} {{\longrightarrow}} \\ & \longrightarrow \mathrm{H} ^ {1} (\mathfrak {g}, \mathrm{L}) \stackrel {{\iota ^ {1}}} {{\longrightarrow}} \mathrm{H} ^ {1} (\mathfrak {g}, \mathrm{M}) \stackrel {{p ^ {1}}} {{\longrightarrow}} \mathrm{H} ^ {1} (\mathfrak {g}, \mathrm{N}) \stackrel {{\delta^ {1}}} {{\longrightarrow}} \dots \end{array}
\]

是一个正合序列。

\begin{enumerate}
\def\labelenumi{(\alph{enumi})}
\setcounter{enumi}{6}
\tightlist
\item
  按 (f) 的记号，正合序列：
\end{enumerate}

\[
0 \to \mathscr {L} (\mathrm{N}, \mathrm{L}) \to \mathscr {L} (\mathrm{N}, \mathrm{M}) \to \mathscr {L} (\mathrm{N}, \mathrm{N}) \to 0
\]

定义了一个正合序列

\[
\mathrm{H} ^ {0} (\mathfrak {g}, \mathscr {L} (\mathrm{N}, \mathrm{M})) \rightarrow \mathrm{H} ^ {0} (\mathfrak {g}, \mathscr {L} (\mathrm{N}, \mathrm{N})) \rightarrow \mathrm{H} ^ {1} (\mathfrak {g}, \mathscr {L} (\mathrm{N}, \mathrm{L})).
\]

 N 的恒等映射是\(\mathcal{L}(\mathrm{N},\mathrm{N})\)中的一个不变元，因而是\(\mathrm{H}^{0}(\mathfrak{g},\mathcal{L}(\mathrm{N},\mathrm{N}))\)中的一个元素；令 c 为它在\(\mathrm{H}^{1}(\mathfrak{g},\mathcal{L}(\mathrm{N},\mathrm{L}))\)中的像。于是，M 中存在补充 L 的子 g-模的充要条件是 c=0。（此条件意味着 u 是\(\mathrm{H}^{0}(\mathfrak{g},\mathcal{L}(\mathrm{N},\mathrm{M}))\)中某个元素的像，即从 g-模 N 到 g-模 M 的一个同态；此时\(v(\mathrm{N})\)就是所求的补子。）

\begin{enumerate}
\def\labelenumi{(\alph{enumi})}
\setcounter{enumi}{7}
\item
  证明\(Z^{1}(\mathfrak{g},\mathfrak{g})\)（其中 g 通过伴随表示视为 g-模）可与 g 的导子空间等同，且\(B^{1}(\mathfrak{g},\mathfrak{g})\)可与 g 的内导子空间等同。
\item
  令\(\mathfrak{a}\)和\(\mathfrak{b}\)为两个李 K-代数，其中\(\mathfrak{b}\)为交换代数，且\(\mathfrak{b} \xrightarrow{\lambda} \mathfrak{g} \xrightarrow{\mu} \mathfrak{a}\)是\(\mathfrak{a}\)通过\(\mathfrak{b}\)的扩张。对于所有\(x \in \mathfrak{g}\)，\(\mathrm{ad}_{\mathfrak{a}} x\)在\(\mathfrak{b}\)上的限制只取决于\(x\)模\(\mathfrak{b}\)的类，因此\(\mathfrak{a}\)上具有\(\mathfrak{b}\)-模结构。令\(\nu\)为从\(\mathfrak{a}\)到\(\mathfrak{g}\)的 K-线性映射，满足\(\mu \circ \nu\)是\(\mathfrak{a}\)的恒等映射。对于\(x, y\)中的\(\mathfrak{a}\)，记\(f(x, y) = [\nu x, \nu y] - \nu([x, y])\)。证明\(f \in Z^2(\mathfrak{a}, \mathfrak{b})\)，且\(c\)（即\(f\)在\(H^2(\mathfrak{a}, \mathfrak{b})\)中的类）不依赖于\(\nu\)的选择。对于\(\mathfrak{a}\)通过\(\mathfrak{b}\)的两个扩张，若它们定义了相同的\(\mathfrak{a}\)-模结构（作用于\(\mathfrak{b}\)），则二者等价的充要条件是相应的类\(c \in H^2(\mathfrak{a}, \mathfrak{b})\)相同。该扩张无本质的充要条件是\(c = 0\)。若 B 是\(\mathfrak{a}\)-模且\(c \in H^2(\mathfrak{a}, \mathfrak{B})\)，则存在\(\mathfrak{a}\)通过 B（将 B 视为交换李代数）的扩张，它在 B 上定义给定的\(\mathfrak{a}\)-模结构，并实现\(c\)这一\(H^2(\mathfrak{a}, \mathfrak{B})\)中的给定元素。
\item
  令 g 为定义在 K 上的有限维李代数，U 为其包络代数，h 为 g 的理想，β 为定义在 g 上且在 h 上限制非退化的不变双线性形式，\((y_{i})_{1 \leqslant i \leqslant n}\)和\((y_{i}')_{1 \leqslant i \leqslant n}\)为 h 的两组基，满足
\end{enumerate}

 \(\beta(y_i, y_i') = \delta_{ij}, t = \sum_{i=1}^{n} y_i y_i' \in \mathbf{U}, \mathbf{M}\)，\(g\)-模中的\(\rho\)是\(\mathrm{C}^*(\mathfrak{g}, \mathbf{M})\)的自同态，它将每个\(\mathrm{C}^p(\mathfrak{g}, \mathbf{M})\)映到\(\mathrm{C}^{p-1}(\mathfrak{g}, \mathbf{M})\)，并且当\(p > 0\)时由下式定义：

\[
(\wp f) (x _ {1}, \dots , x _ {p - 1}) = \sum_ {k = 1} ^ {n} (y _ {k}) _ {\mathrm{M}} f (y _ {k} ^ {\prime}, x _ {1}, \dots , x _ {p - 1}).
\]

最后，令\(\Gamma\)为\(C^*(g, M)\)的自同态，它延拓\(t_M\)，将每个上链\(f\)（次数为\(>0\)）映到上链\(t_M \circ f\)。证明\(\rho d + d\rho = \Gamma\)。（若对\(x \in g\)，写作：

\[
[ x, y _ {k} ] = \sum_ {l = 1} ^ {n} a _ {k l} y _ {l}, \quad [ x, y _ {k} ^ {\prime} ] = \sum_ {l = 1} ^ {n} a _ {k l} ^ {\prime} y _ {l} ^ {\prime},
\]

·先证明\(a_{kl} = -a_{lk}^{\prime}\)。)由此推出\(\Gamma d = d\Gamma\)。若 M 是 K 上有限维的简单模，\(\beta\)是相应的双线性形式，且 dim\(\mathfrak{h}\)不被 K 的特征整除，证明对所有 p 都有\(\mathrm{H}^{p}(\mathfrak{g},\mathrm{M})=\{0\}\)。~（利用命题 12，证明\(\Gamma\)是\(\mathrm{C}^{*}(\mathfrak{g},\mathrm{M})\)的自同构，因而在\(\mathrm{Z}^{p}(\mathfrak{g},\mathrm{M})\)和\(\mathrm{B}^{p}(\mathfrak{g},\mathrm{M})\)上诱导自同构。对于

\[
f \in \mathrm{Z} ^ {p} (\mathfrak {g}, \mathrm{M}), \quad \Gamma f = (d \rho + \rho d) f = d \rho f \in \mathrm{B} ^ {p} (\mathfrak {g}, \mathrm{M}),
\]

于是\(Z^{p}(\mathfrak{g},\mathbf{M}) = \mathrm{B}^{p}(\mathfrak{g},\mathbf{M}).\)

\exercisesfor{4}

第 4 节的约定除非另有说明均有效。

\begin{enumerate}
\def\labelenumi{\arabic{enumi}.}
\item
  令\(\mathfrak{g}\)为幂零李代数，\(p\)（分别地，\(q\)）是满足\(\mathcal{C}^p\mathfrak{g} = \{0\}\)（分别地\(\mathcal{C}_q\mathfrak{g} = \mathfrak{g}\)）的最小整数。证明\(p = q + 1\)，且\(\mathcal{C}_i\mathfrak{g} \supset \mathcal{C}^{p - i}\mathfrak{g}\)。（使用命题 1 的论证。）
\item
  令\(\mathfrak{g}\)为维数为 1 的代数\(\mathfrak{h}\)与交换理想\(\mathfrak{g}'\)的半直积。令\(x \in \mathfrak{h}, x \neq 0\)，且\(u\)为\(\operatorname{ad}_{\mathfrak{g}} x\)在\(\mathfrak{g}'\)上的限制。
\end{enumerate}

\begin{enumerate}
\def\labelenumi{(\alph{enumi})}
\item
  g 为幂零的充要条件是 u 为幂零。
\item
  要使\(\mathfrak{g}\)的 Killing 形式为零，充要条件是\(\operatorname{Tr}(u^2) = 0\)。
\item
  由 (a) 和 (b) 推出，存在 Killing 形式为零的非幂零李代数。
\item
  由 (a) 推出，在满足\(\mathcal{C}^{p-1}\mathfrak{g} \neq \{0\}\)且\(\mathcal{C}^p\mathfrak{g} = \{0\}\)的幂零李代数中，可能有\(\mathcal{C}_i\mathfrak{g} \neq \mathcal{C}^{p-i}\mathfrak{g}\)。
\end{enumerate}

\begin{enumerate}
\def\labelenumi{\arabic{enumi}.}
\setcounter{enumi}{2}
\tightlist
\item
  \begin{enumerate}
  \def\labelenumii{(\alph{enumii})}
  \tightlist
  \item
    令\(\mathfrak{g}\)为幂零李代数，\(\mathfrak{z}\)为其中心，\(\mathfrak{b}\)为\(\mathfrak{g}\)的非零理想。证明\(\mathfrak{z} \cap \mathfrak{b} \neq \{0\}\)。（将\(\mathfrak{b}\)通过伴随表示视为\(\mathfrak{g}\)-模。）
  \end{enumerate}
\end{enumerate}

\begin{enumerate}
\def\labelenumi{(\alph{enumi})}
\setcounter{enumi}{1}
\tightlist
\item
  若在李代数\(\mathfrak{g}\)中，理想\(\mathfrak{b}\)包含于\(\mathcal{C}_{i+1}\mathfrak{g}\)但不包含于\(\mathcal{C}_i\mathfrak{g}\)，证明理想\(\mathfrak{b} \cap \mathcal{C}_k\mathfrak{g}\)在\(0 \leqslant k \leqslant i + 1\)所给范围内两两不同。（应用 (a)。）
\end{enumerate}

\begin{enumerate}
\def\labelenumi{\arabic{enumi}.}
\setcounter{enumi}{3}
\tightlist
\item
  令 g 为幂零李代数。
\end{enumerate}

\begin{enumerate}
\def\labelenumi{(\alph{enumi})}
\item
  \(\mathfrak{g}\)的每个子代数都是次不变的。（使用命题 3。）
\item
  令\(\mathfrak{b}\)为\(\mathfrak{g}\)的向量子空间，满足\(\mathfrak{b} + \mathcal{D}\mathfrak{g} = \mathfrak{g}\)。证明\(\mathfrak{g}\)的子代数由\(\mathfrak{b}\)生成，且等于\(\mathfrak{g}\)。（对该子代数应用 (a)，从而将问题化为\(\mathfrak{b}\)是\(\mathfrak{g}\)的理想的情形，并使用 § 1 的命题 4。）推出\(\mathfrak{g}\)的最少生成元数为\(\dim \mathfrak{g}/\mathcal{D}\mathfrak{g}\)。
\end{enumerate}

\begin{enumerate}
\def\labelenumi{\arabic{enumi}.}
\setcounter{enumi}{4}
\tightlist
\item
  \begin{enumerate}
  \def\labelenumii{(\alph{enumii})}
  \tightlist
  \item
    证明：若李代数 g 的每个子代数都是次不变的，则它是幂零的。（证明：若\(\dim g > 1\)，则每个\(x \in g\)都属于 g 的一个理想\(h \neq g\)，该理想具有与 g 相同的性质，因而由关于 g 维数的归纳假设可知它是幂零的；所以\(ad_{g} x\)是幂零的。）
  \end{enumerate}
\end{enumerate}

\begin{enumerate}
\def\labelenumi{(\alph{enumi})}
\setcounter{enumi}{1}
\tightlist
\item
  证明：若在李代数\(\mathfrak{g}\)中，每个不同于\(\mathfrak{g}\)的子代数都不同于其正规化子，则\(\mathfrak{g}\)是幂零的。（将其化归为 (a)。）
\end{enumerate}

\begin{enumerate}
\def\labelenumi{\arabic{enumi}.}
\setcounter{enumi}{5}
\item
  令 g 为幂零李代数，a 为 g 的交换理想。以下条件等价：(a) a 是 g 的极大交换理想；(b) a 是 g 的极大交换子代数；(c) a 等于其在 g 中的中心化子\(a'\)。（非 (a)\(\Rightarrow\)非 (b)\(\Rightarrow\)非 (c) 的蕴含显然。若\(a' \neq a\)，由命题 1 存在 g 的理想\(a''\)，满足\(a \subset a'' \subset a'\)且 dim\(a''/a = 1\)。于是\(a'' = a + Kx\)，从而\([a'', a''] \subset [a, a] + [x, a] = \{0\}\)，故 (a) 不成立。）
\item
  \begin{enumerate}
  \def\labelenumii{(\alph{enumii})}
  \tightlist
  \item
    令 V 为定义在\(K\)上的有限维向量空间，\(\mathfrak{g}\)为 V 上幂零自同态组成的李代数，\((\mathrm{V}_{r})_{0 \leq r \leq n}\)为 V 的递减向量子空间序列，满足\(V_{0} = V, V_{n} = \{0\}\)，且\(\mathfrak{g}(\mathrm{V}_{r}) \subset V_{r+1}\)在\(0 \leq r < n\)时成立。通过对 i 归纳证明\((\mathcal{D}^{i}\mathfrak{g})(\mathrm{V}_{r}) \subset V_{r+2^{i}}\)。若\(\dim V \geq 2^{i}\)，则 V 中被\(\mathcal{D}^{i}\mathfrak{g}\)湮灭的元素所成子空间的维数\(\geqslant 2^{i}\)。若\(\dim V \leqslant 2^{i}\)，则\(\mathcal{D}^{i}\mathfrak{g} = \{0\}\)。
  \end{enumerate}
\end{enumerate}

\begin{enumerate}
\def\labelenumi{(\alph{enumi})}
\setcounter{enumi}{1}
\item
  令\(\mathfrak{g}\)为幂零李代数，\(i\)为满足\(\geqslant 0\)的整数。若\(\dim \mathcal{D}^{i}\mathfrak{g} > 2^{i} + 1\)，则\(\mathcal{D}^{i}\mathfrak{g}\)的中心维数\(\geqslant 2^{i}\)。若\(\dim \mathcal{D}^{i}\mathfrak{g} \leqslant 2^{i} + 1\)，则\(\mathcal{D}^{i}\mathfrak{g}\)为交换的。（将 (a) 应用于\(\mathcal{D}^{i}\mathfrak{g}\)上的\(\operatorname{ad}_{\mathfrak{g}} x, x \in \mathfrak{g}\)的限制。）
\item
  令\(\mathfrak{g}\)为幂零李代数，\(i\)为满足\(\geqslant 0\)的整数。若\(\mathcal{D}^{i}\mathfrak{g}\)不交换，则\(\mathcal{D}^{i}\mathfrak{g}/\mathcal{D}^{i+1}\mathfrak{g}\)的维数\(\geqslant 2^{i} + 1\)。（将其转到商代数，化归为\(\dim \mathcal{D}^{i+1}\mathfrak{g} = 1\)的情形，并使用 (b)。）
\end{enumerate}

\begin{enumerate}
\def\labelenumi{\arabic{enumi}.}
\setcounter{enumi}{7}
\tightlist
\item
  令\(\mathfrak{g}\)为李代数，\(\mathfrak{m}\)为余维 1 的理想，\(x\)为\(\mathfrak{g}\)中不属于\(\mathfrak{m}\)的元素，且\(z \in \mathfrak{g}\)。
\end{enumerate}

\begin{enumerate}
\def\labelenumi{(\alph{enumi})}
\item
  证明线性映射 D 在 m 上为 0 并将 x 映到 z 时，若 z 属于 g 中 m 的中心化子 a，则 D 是导子。
\item
  令\(q\)为满足\(a \subset \mathcal{C}^q\mathfrak{g}\)的最大整数。证明若进一步有\(z \notin \mathcal{C}^{q+1}\mathfrak{g}\)，则\(D\)不是\(\mathfrak{g}\)的内导子。
\item
  由 (a)、(b) 和习题 7 (c) 推出：若 g 是维数为\textgreater1 的幂零李代数，则 g 的内导子空间在全体导子空间中的余维数为\(\geqslant2\)。
\end{enumerate}

\begin{enumerate}
\def\labelenumi{\arabic{enumi}.}
\setcounter{enumi}{8}
\tightlist
\item
  \begin{enumerate}
  \def\labelenumii{(\alph{enumii})}
  \tightlist
  \item
    验证下列乘法表定义了维数分别为 3 和 4 的两个幂零李代数\(\mathfrak{g}_3, \mathfrak{g}_4\)：
  \end{enumerate}
\end{enumerate}

\[
\mathfrak {g} _ {3}: \left[ x _ {1}, x _ {2} \right] = x _ {3}, \quad \left[ x _ {1}, x _ {3} \right] = \left[ x _ {2}, x _ {3} \right] = 0
\]

\[
\mathfrak {g} _ {4}: \left[ x _ {1}, x _ {2} \right] = x _ {3}, \quad \left[ x _ {1}, x _ {3} \right] = x _ {4}, \quad \left[ x _ {1}, x _ {4} \right] = \left[ x _ {2}, x _ {3} \right] = \left[ x _ {2}, x _ {4} \right] = \left[ x _ {3}, x _ {4} \right] = 0.
\]

\begin{enumerate}
\def\labelenumi{(\alph{enumi})}
\setcounter{enumi}{1}
\item
  证明\(\mathfrak{g}_3\)同构于\(\mathfrak{n}(3, \mathbf{K})\)，也同构于\(\mathfrak{sl}(2, \mathbf{K})\)；当\(\mathbf{K}\)的特征为 2 时上述结论成立。
\item
  令\(\mathfrak{g}_1\)为唯一的 1 维李代数。证明维数\(\leqslant 4\)的幂零李代数如下表所示：
\end{enumerate}

\[
\text { dimension } 1: \mathfrak {g} _ {1}; \text { dimension } 2: (\mathfrak {g} _ {1}) ^ {2};
\]

维数 3: \((\mathfrak{g}_1)^3, \mathfrak{g}_3\) ; 维数 4: \((\mathfrak{g}_1)^4, \mathfrak{g}_3 \times \mathfrak{g}_1, \mathfrak{g}_4\) .

（使用习题 7 (c)。若\(\mathfrak{g}\)是幂零的且维数为 3，并且\(\dim \mathcal{D}\mathfrak{g} = 1\)，注意到\(\mathcal{D}\mathfrak{g}\)包含于\(\mathfrak{g}\)的中心，从而\(\mathfrak{g} = \mathfrak{g}_3\)。若 dim \(\mathfrak{g} = 4\)，dim \(\mathcal{D}\mathfrak{g} = 1\)，注意括号运算在\(\mathfrak{g}\)上定义了\(\mathfrak{g} / \mathcal{D}\mathfrak{g}\)上的交错双线性形式；该形式必然退化，因此存在子代数\(\mathfrak{h}\)，它是\(\mathfrak{g}\)的子代数且包含于\(\mathfrak{g}\)的中心，满足 dim \(\mathfrak{h} = 1\)且\(\mathfrak{h}\cap \mathcal{D}\mathfrak{g} = \{0\}\)，从而\(\mathfrak{g} = \mathfrak{h}\times \mathfrak{g}_3\)。若 dim \(\mathfrak{g} = 4\)，dim \(\mathcal{D}\mathfrak{g} = 2\)，则\(\mathcal{D}\mathfrak{g}\)为交换的；将定理 1 应用于\(\mathcal{D}\mathfrak{g}\)上\(\operatorname{ad}_{\mathfrak{g}} x\)的限制\((x\in \mathfrak{g})\)，证明存在\(\mathfrak{h}\)的交换理想\(\mathfrak{g}\)，满足\(\mathfrak{h}\supseteq \mathcal{D}\mathfrak{g}\)且 dim \(\mathfrak{h} = 3\)；令\(x\in \mathfrak{g},\)，且\(x\notin \mathfrak{h}\)；选取\(\mathfrak{h}\)的一组基，使\(\operatorname{ad}_{\mathfrak{g}}x\)在\(\mathfrak{h}\)上的限制关于此基具有 Jordan 矩阵；于是\(\mathfrak{g} = \mathfrak{g}_4\)。

\begin{enumerate}
\def\labelenumi{\arabic{enumi}.}
\setcounter{enumi}{9}
\tightlist
\item
  令\(\mathcal{D}\)为\(\mathfrak{g}_4\)的导子李代数。证明\(\mathcal{D}\)的维数为 7、中心为零，并且内导子理想\(\mathcal{D}' \subset \mathcal{D}\)在\(\mathcal{D}\)中没有补子代数。
\end{enumerate}

¶ 11. 令 L 为交换环，A 为 L 上的左阿廷代数，γ 为从 A × A 到 L 的映射。在 A 上定义内部复合律\((a, b) \mapsto a * b = ab + \gamma(a, b)ba\)。对于 A 的每个子集 E，令\(\tilde{E}\)表示由 E 生成的 A 的子环（无单位元）。证明若 E 由幂零元组成且关于运算 * 稳定，则\(\tilde{E}\)是幂零的（即存在 n\textgreater{}0，使得\(\tilde{E}^{n} = \{0\}\)）。按以下步骤进行：

\begin{enumerate}
\def\labelenumi{\arabic{enumi}.}
\item
  首先设 A 是单环，因此同构于\(\mathcal{L}_{\mathrm{D}}(\mathrm{T})\)，其中 T 是域 D 上维数为\(m\)的左向量空间。对\(m\)作归纳。令\(\Phi\)为由满足以下条件的子集组成的集合：\(\mathrm{F} \subset \mathrm{E}\)关于\(*\)稳定，且\(\tilde{\mathrm{F}}\)是幂零的。证明\(\tilde{\Phi}\)具有极大元 M（注意，\(\tilde{\mathrm{F}}^m = \{0\}\)对所有\(\mathrm{F} \in \Phi\)都成立）。假设\(\tilde{\mathbf{M}} \neq \tilde{\mathbf{E}}\)。证明存在\(a \in \mathrm{E}\)，使得\(a \notin \tilde{\mathbf{M}}\)，且\(t * a \in \tilde{\mathbf{M}}\)对所有\(t \in \mathbf{M}\)都成立（注意，不存在无限序列\((a_n)\)，使得\(a_n \in \mathrm{E}\)、\(a_n \notin \mathbf{M}\)、\(a_n = t_{n-1} * a_{n-1}\)，且\(t_{n-1} \in \mathbf{M}\)）。令 S 为 T 的子空间，它是各个\(u(\mathrm{T})\)的直和，其中\(u \in \tilde{\mathbf{M}}\)；证明\(S \neq \{0\}\)且\(S \neq T\)，并且\(a(S) \subset S\)。令 N 为所有\(v \in E\)中满足\(v(S) \subset S\)者的集合。利用归纳假设，并把 N 的元素视为在 S 及 T/S 上的作用，证明\(N \in \Phi\)，从而矛盾。
\item
  一般情形下，利用 A 的 Jacobson 根基是幂零的这一事实。由此结果得到 Engel 定理的一个新证明。
\item
  令\(\mathfrak{g}\)为李代数，\(\mathcal{C}^{\infty}\mathfrak{g}\)为各个\(\mathcal{C}^p\mathfrak{g}\)的交。
\end{enumerate}

\begin{enumerate}
\def\labelenumi{(\alph{enumi})}
\item
 李代数\(\mathfrak{g} / \mathcal{C}^{\infty}\mathfrak{g}\)是幂零的。
\item
  证明存在幂零子代数\(\mathfrak{h}\)，它是\(\mathfrak{g}\)的子代数，满足\(\mathfrak{g} = \mathfrak{h} + \mathcal{C}^{\infty}\mathfrak{g}\)。（对\(\dim \mathfrak{g}\)作归纳。假设\(\mathfrak{g}\)不是幂零的。取\(x\in \mathfrak{g}\)，使得\(\mathbf{I} = \bigcap_{k}\left(\mathrm{(ad}x)^{k}(\mathfrak{g})\neq \{0\}\right)\)。令\(\mathfrak{n}\)为\((\mathrm{ad}x)^k\)的核的并。证明这是\(\mathfrak{g}\)的子代数，且\(\mathfrak{g}\)是\(\mathbf{I}\)与\(\mathfrak{n}\)的直和。由归纳假设，\(\mathfrak{n}\)是\(\mathcal{C}^{\infty}\mathfrak{n}\)与幂零子代数\(\mathfrak{h}\)之和。最后，\(\mathcal{C}^{\infty}\mathfrak{n}\subset \mathcal{C}^{\infty}\mathfrak{g}\)且\(\mathbf{I}\subset \mathcal{C}^{\infty}\mathfrak{g}\)。
\item
  验证下列乘法表定义了一个 5 维（可解）李代数 g：
\end{enumerate}

\[
\left[ x _ {1}, x _ {2} \right] = x _ {5}, \quad \left[ x _ {1}, x _ {3} \right] = x _ {3}, \quad \left[ x _ {1}, x _ {4} \right] = - x _ {4}, \quad \left[ x _ {3}, x _ {4} \right] = x _ {5}
\]

且其他\([x_i, x_j]\)均为零。证明对于此李代数

\[
\mathcal {C} ^ {\infty} \mathfrak {g} = \mathrm{K} x _ {3} + \mathrm{K} x _ {4} + \mathrm{K} x _ {5}
\]

但\(\mathcal{C}^{\infty}\mathfrak{g}\)在\(\mathfrak{g}\)中不存在补子代数。

¶ 13. 令 g 为李代数，\(\mathfrak{z}\)为 g 中\(C^{\infty}g\)的中心化子。若\(\mathfrak{z} \not\subset C^{\infty}g\)，则 g 的中心非零。（写成\(g = C^{\infty}g + t\)，其中 t 是 g 的幂零子代数（习题 12）。令\(g_{1} = \mathfrak{z} + t\)。令 b 为幂零子代数，满足\(g_{1} = C^{\infty}g_{1} + \mathfrak{b}\)。证明\(g = C^{\infty}g + \mathfrak{b}\)且\(C^{\infty}g_{1} \subset \mathfrak{z} \cap C^{\infty}g\)。令 y 为\(\mathfrak{z}\)中不属于\(C^{\infty}g\)的元素。写成\(y = x + x'\)，其中\(x \in \mathfrak{b}, x' \in C^{\infty}g_{1}\)；于是\(x \in \mathfrak{z}\)且\(x \neq 0\)，因此\(\mathfrak{z} \cap \mathfrak{b}\)是 b 的非零理想；利用习题 3 (a)，推出 g 中存在非零元素，它与\(C^{\infty}g\)及 b 可交换。）

¶ 14. 令 g 为李代数，b 为 g 的次不变子代数，且 b 在 g 中的中心化子\(\mathfrak{z}(b)\)为零。

\begin{enumerate}
\def\labelenumi{(\alph{enumi})}
\item
  中心化子\(\mathfrak{z}(\mathcal{C}^{\infty}\mathfrak{b})\)（即\(\mathcal{C}^{\infty}\mathfrak{b}\)在\(\mathfrak{g}\)中的中心化子）包含于\(\mathcal{C}^{\infty}\mathfrak{b}\)。（若\(\mathfrak{z}(\mathcal{C}^{\infty}\mathfrak{b})\not\subset\mathcal{C}^{\infty}\mathfrak{b},\mathfrak{z}(\mathcal{C}^{\infty}\mathfrak{b})\not\subset\mathfrak{b}\)，由习题 13 得；\(\mathcal{C}^{\infty}\mathfrak{b}\)是\(\mathfrak{g}\)的理想（§ 1，习题 14）；令\(\mathfrak{c}=\mathfrak{b}+\mathfrak{z}(\mathcal{C}^{\infty}\mathfrak{b})\)；\(\mathfrak{c}\)是子代数，\(\mathfrak{b}\)在\(\mathfrak{c}\)中是次不变的；\(\mathfrak{b}\)是\(\mathfrak{b}_{1}\subset\mathfrak{c}\)的理想，其中\(\mathfrak{b}\neq\mathfrak{b}_{1}\)；令\(y\in\mathfrak{b}_{1},y\notin\mathfrak{b};y=x+x'\)，其中\(x\in\mathfrak{z}(\mathcal{C}^{\infty}\mathfrak{b}),x'\in\mathfrak{b}\)；则\(x\in\mathfrak{z}(\mathcal{C}^{\infty}\mathfrak{b})\cap\mathfrak{b}_{1},x\notin\mathfrak{b};\mathfrak{b}_{1}=\mathrm{Kx}+\mathfrak{b}\)是子代数；\(\mathcal{C}^{\infty}\mathfrak{c}=\mathcal{C}^{\infty}\mathfrak{b}\)，因为\(\mathfrak{c}\cap\mathfrak{z}(\mathcal{C}^{\infty}\mathfrak{b})\subset\mathcal{C}^{\infty}\mathfrak{b}\)，从而\(\mathcal{C}^{k}\mathfrak{c}\subset\mathcal{C}^{k}\mathfrak{b}\)对所有\(k\)均成立；\(\mathfrak{z}(\mathcal{C}^{\infty}\mathfrak{b})\cap\mathfrak{b}\not\subset\mathcal{C}^{\infty}\mathfrak{b}\)，因此\(\mathfrak{b}\)的中心由习题 13 知非零；特别地，\(\mathfrak{z}(\mathfrak{b})\neq\{0\}\)，矛盾。）
\item
  由 (a) 推出，若\(\mathfrak{D}\)表示\(\mathcal{C}^{\infty}\mathfrak{b}\)的导子李代数，\(\mathfrak{a}\)表示\(\mathcal{C}^{\infty}\mathfrak{b}\)的中心，则\(\dim \mathfrak{g} \leqslant \dim \mathfrak{D} + \dim \mathfrak{a}\)（令\(\mathfrak{g}\)通过伴随表示作用于\(\mathcal{C}^{\infty}\mathfrak{b}\)。）
\end{enumerate}

\begin{enumerate}
\def\labelenumi{\arabic{enumi}.}
\setcounter{enumi}{14}
\tightlist
\item
  令\(l\)为中心为零的李代数，\(\mathfrak{D}_1\)为\(l\)的导子李代数，\(\mathfrak{D}_2\)为\(\mathfrak{D}_1, \ldots\)的导子李代数。于是\(l\)是\(\mathfrak{D}_1, \mathfrak{D}_1\)的理想，并且理想链继续到\(\mathfrak{D}_2\)，等等（§ 1，习题 15 (c)）。令\(\mathfrak{D}\)为\(\mathcal{C}^\infty l\)的导子李代数，\(a\)为\(\mathcal{C}^\infty l\)的中心。
\end{enumerate}

\begin{enumerate}
\def\labelenumi{(\alph{enumi})}
\item
  于是\(\dim \mathfrak{D}_1 \leqslant \dim \mathfrak{D} + \dim \mathfrak{a}\)。（使用上面的习题 14 和 § 1 的习题 15。）
\item
  由 (a) 推出，当 i 充分大时，\(D_{i}\)的所有导子都是内导子。
\end{enumerate}

\begin{enumerate}
\def\labelenumi{\arabic{enumi}.}
\setcounter{enumi}{15}
\item
  令\(\mathfrak{g}_1\)和\(\mathfrak{g}_2\)为李 K-代数，\(\mathfrak{n}_1\)和\(\mathfrak{n}_2\)分别为它们的最大幂零理想。证明\(\mathfrak{g}_1 \times \mathfrak{g}_2\)的最大幂零理想是\(\mathfrak{n}_1 \times \mathfrak{n}_2\)。
\item
  我们重新考察习题 9 (a) 中的李代数\(\mathfrak{g}_3\)，其基为\((x_1, x_2, x_3)\)。
\end{enumerate}

\begin{enumerate}
\def\labelenumi{(\alph{enumi})}
\tightlist
\item
  令 V 为向量空间 K{[}X{]}。令 D 为 V 上关于 X 的微分算子，M 为乘以 X 的算子。证明若 K 的特征为 0，则映射
\end{enumerate}

\[
\alpha x _ {1} + \beta x _ {2} + \gamma x _ {3} \mapsto \alpha D + \beta M + \gamma
\]

是 V 上的无限维简单表示\(\rho\)，其表示的李代数为\(\mathfrak{g}\)。

\begin{enumerate}
\def\labelenumi{(\alph{enumi})}
\setcounter{enumi}{1}
\tightlist
\item
  若 K 的特征为 p\textgreater{}0，则理想\((\mathbf{X}^{p})\)（即 K{[}X{]}的理想）在\(\rho(\mathfrak{g})\)下稳定。在 g 在\(\mathrm{K}[\mathrm{X}] / (\mathrm{X}^{p})\)上的商表示下，没有直线稳定。
\end{enumerate}

\begin{enumerate}
\def\labelenumi{\arabic{enumi}.}
\setcounter{enumi}{17}
\tightlist
\item
  令\(\mathfrak{g}\)为定义在\(\mathbf{K}\)上的 7 维向量空间，基为\((e_i)_{1\leqslant i\leqslant 7}\)。在\(\mathfrak{g}\)上按下列公式定义交错括号：
\end{enumerate}

\[
\left[ e _ {i}, e _ {j} \right] = \alpha_ {i j} e _ {i + j} \quad (1 \leqslant i <   j \leqslant 7, i + j \leqslant 7)\tag{1}
\]

且其他括号\([e_i, e_j]\)在\(i < j\)时均为零。

\begin{enumerate}
\def\labelenumi{(\alph{enumi})}
\tightlist
\item
  要使 Jacobi 恒等式成立，充要条件是：
\end{enumerate}

\begin{enumerate}
\def\labelenumi{(\arabic{enumi})}
\setcounter{enumi}{1}
\tightlist
\item
\end{enumerate}

\[
- \alpha_ {2 3} \alpha_ {1 5} + \alpha_ {1 3} \alpha_ {2 4} = 0\tag{2}
\]

\[
\alpha_ {1 2} \alpha_ {3 4} - \alpha_ {2 4} \alpha_ {1 6} + \alpha_ {1 4} \alpha_ {2 5} = 0.\tag{3}
\]

\begin{enumerate}
\def\labelenumi{(\alph{enumi})}
\setcounter{enumi}{1}
\item
  以下设所有\(\alpha_{ij}\)均\(\neq 0\)。证明理想\(\mathcal{C}^2\mathfrak{g}\)、\(\mathcal{C}^3\mathfrak{g}\)、\(\mathcal{C}^4\mathfrak{g}\)、\(\mathcal{C}^5\mathfrak{g}\)、\(\mathcal{C}^6\mathfrak{g}\)分别具有以下基：\((e_3, e_4, e_5, e_6, e_7)\)、\((e_4, e_5, e_6, e_7)\)、\((e_5, e_6, e_7)\)、\((e_6, e_7)\)、\((e_7)\)。证明\(\mathfrak{h}\)的中心化子\(\mathcal{C}^5\mathfrak{g}\)具有基\((e_2, e_3, e_4, e_5, e_6, e_7)\)。
\item
  令\((e_i')_{1 \leqslant i \leqslant 7}\)为\(\mathfrak{g}\)的另一组基，满足\(\mathcal{C}^6\mathfrak{g}\)具有基\((e_7')\)，\(\mathcal{C}^5\mathfrak{g}\)具有基\((e_6', e_7')\)，\ldots\(\mathfrak{h}\)具有基\((e_2', e_3', e_4', e_5', e_6', e_7')\)。证明
\end{enumerate}

\[
\left[ e _ {i} ^ {\prime}, e _ {j} ^ {\prime} \right] = \alpha_ {i j} ^ {\prime} e _ {i + j} ^ {\prime} + \beta e _ {i + j + 1} ^ {\prime} + \gamma e _ {i + j + 2} ^ {\prime} + \dots + \lambda e _ {7} ^ {\prime},
\]

其中

\[
\alpha_ {1 4} ^ {\prime} \alpha_ {2 5} ^ {\prime} \alpha_ {1 6} ^ {\prime - 1} \alpha_ {2 4} ^ {\prime - 1} = \alpha_ {1 4} \alpha_ {2 5} \alpha_ {1 6} ^ {- 1} \alpha_ {2 4} ^ {- 1}.
\]

\begin{enumerate}
\def\labelenumi{(\alph{enumi})}
\setcounter{enumi}{3}
\tightlist
\item
  由此推出：若 K 是无限域，则存在无穷多个互不同构的 7 维 K 上幂零李代数；并且存在不能通过将基域从 R 扩张到 C 而由实幂零李代数得到的复幂零李代数。
\end{enumerate}

\begin{enumerate}
\def\labelenumi{\arabic{enumi}.}
\setcounter{enumi}{18}
\tightlist
\item
  \begin{enumerate}
  \def\labelenumii{(\alph{enumii})}
  \tightlist
  \item
    令\(\mathfrak{g}\)为李代数。令\(\mathfrak{S}\)为\(\mathfrak{g}\)的导子李代数。\(\mathfrak{g}^{[k]}\)的特征理想（在\(\mathfrak{g}\)中）按如下归纳方式定义：\(\mathfrak{g}^{[0]} = \mathfrak{g}\)，且\(\mathfrak{g}^{[k+1]}\)是由\(\mathfrak{g}\)中的 Dx 生成的子空间（\(\mathrm{D} \in \mathfrak{S}, x \in \mathfrak{g}^{[k]}\)）。以下条件等价：(1)\(\mathfrak{g}\)的每个导子都是幂零的；(2)当\(\mathfrak{g}^{[k]} = \{0\}\)在\(k\)充分大时成立；(3)\(\mathfrak{g}\)的全纯代数是幂零的。若这些条件成立，则\(\mathfrak{g}\)称为特征幂零。这样的代数是幂零的。
  \end{enumerate}
\end{enumerate}

\begin{enumerate}
\def\labelenumi{(\alph{enumi})}
\setcounter{enumi}{1}
\tightlist
\item
  验证下述乘法表定义了一个 8 维李代数：
\end{enumerate}

\[
\begin{array}{r} \left[ x _ {1}, x _ {2} \right] = x _ {3}, \quad \left[ x _ {1}, x _ {3} \right] = x _ {4}, \quad \left[ x _ {1}, x _ {4} \right] = x _ {5}, \quad \left[ x _ {1}, x _ {5} \right] = x _ {6} \\ \left[ x _ {1}, x _ {6} \right] = x _ {8}, \quad \left[ x _ {1}, x _ {7} \right] = x _ {6}, \quad \left[ x _ {2}, x _ {3} \right] = x _ {5}, \quad \left[ x _ {2}, x _ {4} \right] = x _ {6} \\ \left[ x _ {2}, x _ {5} \right] = x _ {7}, \quad \left[ x _ {2}, x _ {6} \right] = 2 x _ {8}, \quad \left[ x _ {3}, x _ {4} \right] = - x _ {7} + x _ {8}, \quad \left[ x _ {3}, x _ {5} \right] = - x _ {8} \end{array}
\]

且\([x_i, x_j] = 0\)在\(i + j > 8\)时成立；证明

\[
\begin{array}{c c} \mathcal {C} ^ {2} \mathfrak {g} = \sum_ {i = 3} ^ {8} \mathrm{K} x _ {i}, & \mathcal {C} ^ {3} \mathfrak {g} = \sum_ {i = 4} ^ {8} \mathrm{K} x _ {i}, \\ & \mathcal {C} ^ {6} \mathfrak {g} = \mathrm{K} x _ {8}, \quad \mathcal {C} ^ {7} \mathfrak {g} = \{0 \}, \end{array} \quad \begin{array}{c c} \mathcal {C} ^ {4} \mathfrak {g} = \sum_ {i = 5} ^ {8} \mathrm{K} x _ {i}, & \mathcal {C} ^ {5} \mathfrak {g} = \sum_ {i = 6} ^ {8} \mathrm{K} x _ {i}, \\ & [ \mathcal {C} ^ {2} \mathfrak {g}, \mathcal {C} ^ {2} \mathfrak {g} ] = \mathrm{K} x _ {7} + \mathrm{K} x _ {8}, \end{array}
\]

并且\(\mathcal{C}^2\mathfrak{g}\)到\(\mathcal{C}^4\mathfrak{g}\)的输运子空间是\(\sum_{i=2}^{\infty}\mathrm{Kx}_i\)。推出每个\(\mathfrak{g}\)的导子\(\mathrm{D}x_i = \sum_{j=1}^{8}u_{ij}x_j\)均由公式定义。再证明：当 K 的特征不为 2 时，\(u_{ii} = 0\)对所有\(i\)均成立，从而\(\mathfrak{g}\)是特征幂零的。

\begin{enumerate}
\def\labelenumi{(\alph{enumi})}
\setcounter{enumi}{2}
\tightlist
\item
  若 K 的特征为\(\neq2\)，证明 (b) 中的李代数 g 不是任何李代数的导代数。（若 g = Dh，先证明 h 是幂零的。注意\(\dim g/Dg = 2\)；结合习题 7 (c) 得到矛盾。）
\end{enumerate}

¶ 20. 令 g 为李代数，U 为其包络代数。

\begin{enumerate}
\def\labelenumi{(\alph{enumi})}
\item
  令\(\mathfrak{g}'\)为\(\mathfrak{g}\)的理想，\(\mathrm{U}' \subset \mathrm{U}\)为其包络代数，\(x \in \mathfrak{g}\)以及\(a_1, \ldots, a_n\)为 U 中的元素。假设\(a_i = x\)对指标\(j_1, \ldots, j_n\)成立，并且\(a_i \in \mathrm{U}'\)对其他指标成立（令\(k_1, \ldots, k_q\)为这些指标，且满足\(k_1 < k_2 < \cdots < k_q\)）。那么\(a_1 a_2 \ldots a_n - x^p a_{k_1}a_{k_2} \ldots a_{k_q}\)是形如\(x^p b\)的项之和，其中\(b \in \mathrm{U}'\)且\(p' < p\)。（对\(p\)作归纳。）
\item
  此后设 g 为幂零李代数，K 的特征为 0。令 g' 为 g 中余维 1 的理想，U' 为其包络代数，x 为 g 中不属于 g' 的元素，U₀（分别为 U₀'）为 U（分别为 U'）中被 g 的一组导子 b 湮灭的子代数（这些导子延拓为 U 的导子，并将 g 映入 g'）。假设 U₀ 包含于 U 的中心且 U₀ ≠ U'。证明存在 a₁ ∈ U₀'、a₂ ∈ U'，使得 a₁ ≠ 0 且 a = xa₁ + a₂ ∈ U₀。（令 xᵐbₘ + xᵐ⁻¹bₘ₋₁ + ⋯ + b₀ ∈ U₀，其中 bₘ、\ldots、b₀ 属于 U'，且 m\textgreater{}0，bₘ ≠ 0。利用⟨a⟩证明：对 g 的每个导子 D ∈ b，都有 Dbₘ = 0，m(Dx)bₘ + Dbₘ₋₁ = 0，从而 D(mxbₘ + bₘ₋₁) = 0。）证明 U₀ 包含于其分式域中的代数 K{[}a, a₁⁻¹, U₀'{]}，该代数由 a、a₁⁻¹、U₀' 生成（可由 § 2 定理 1 的推论 7 构造）。推出 U₀ 的分式域是由 a 和 U₀' 生成的域。证明 a 在 K(U₀') 上超越。
\item
 令\(\{0\} = \mathfrak{g}_0 \subset \mathfrak{g}_1 \subset \cdots \subset \mathfrak{g}_n = \mathfrak{g}\)为\(\mathfrak{g}\)的维数分别为\(0, 1, \ldots, n\)的理想序列。令\(x_j\)为\(\mathfrak{g}_j\)中不属于\(\mathfrak{g}_{j-1}\)的元素。令\(j_1 < j_2 < \cdots < j_n\)为满足以下条件的指标\(j\)：存在\(\mathrm{U}^j\)（\(\mathfrak{g}_j\)的包络代数）中的\(\mathrm{U}_0\)（\(\mathrm{U}\)的中心），它不属于\(\mathrm{U}^{j-1}\)。根据 (b)，存在\(a_{j_1} \in \mathrm{U}_0 \cap \mathrm{U}^{j_1}\)和\(a_{j_2} \in \mathrm{U}^{j_2-1}\)，满足\(a_{j_1} \neq 0\)且\(a_j = x_j a_{j_1} + a_{j_2} \in \mathrm{U}_0 \cap \mathrm{U}^{j}\)。对\(n\)作归纳证明：
\end{enumerate}

\[
\mathrm{U} _ {0} \subset \mathrm{K} [ a _ {j _ {1}}, \dots , a _ {j _ {q}}, a _ {j _ {1}} ^ {- 1}, \dots , a _ {j _ {q}} ^ {- 1} ]
\]

并且\(U_{0}\)的分式域由代数独立元\(a_{j_{1}},\ldots,a_{j_{q}}\)生成。特别地，U 中心的分式域是 K 的纯超越扩张。

\begin{enumerate}
\def\labelenumi{\arabic{enumi}.}
\setcounter{enumi}{20}
\tightlist
\item
  令 g 为李代数，σ 为 g 的自同构。
\end{enumerate}

\begin{enumerate}
\def\labelenumi{(\alph{enumi})}
\tightlist
\item
  首先设\(\mathbf{K}\)代数闭。对任意\(\lambda \in \mathbf{K}\)，令\(\mathrm{L}_{\lambda}\)为由\(x \in \mathfrak{g}\)中被\(\sigma - \lambda \mathrm{I}\)的幂所湮灭的元素组成的集合；\(\mathfrak{g}\)是各个\(\mathrm{L}_{\lambda}\)的直和。证明\([\mathrm{L}_{\lambda}, \mathrm{L}_{\mu}] \subset \mathrm{L}_{\lambda \mu}\)。（注意
\end{enumerate}

\[
(\sigma - \lambda \mu \mathrm{I}) ([ x, y ]) = [ (\sigma - \lambda \mathrm{I}) x, \sigma y ] + [ \lambda x, (\sigma - \mu \mathrm{I}) y ].)
\]

\begin{enumerate}
\def\labelenumi{(\alph{enumi})}
\setcounter{enumi}{1}
\item
  对任意域 K，设\(\sigma\)的任一特征值（在 K 的代数闭扩张中）都不是单位根。证明 g 是幂零的。（将其化归为 K 代数闭的情形。若\(\lambda_{1},\ldots,\lambda_{m}\)是\(\sigma\)的不同特征值，则\(\lambda_{i}\lambda_{j},\lambda_{i}\lambda_{j}^{2},\lambda_{i}\lambda_{j}^{3},\ldots\)不全是特征值，因此\(ad_{g}x\)对\(x\in L_{\lambda_{i}}\)是幂零的。使用习题 11，将\(ad_{g}x\)应用于其中 x 遍历\(L_{\lambda_{i}}\)的集合。）
\item
  对任意域 K，设\(\sigma^{q}=I\)，其中 q 为素数，且\(\sigma\)没有特征值等于 1。证明 g 是幂零的。（采用 (b) 中的方法，注意：若\(g\neq\{0\}\)，则 q 不等于 K 的特征，并且对\(\lambda_{i},\lambda_{j}\)中的每一对有序特征值\(\sigma\)，存在整数 k 使\(\lambda_{i}\lambda_{j}^{k}=1\)。）
\end{enumerate}

\begin{enumerate}
\def\labelenumi{\arabic{enumi}.}
\setcounter{enumi}{21}
\tightlist
\item
  \begin{enumerate}
  \def\labelenumii{(\alph{enumii})}
  \tightlist
  \item
    令\(u, v\)为有限维向量空间\(\mathbf{E}\)上的两个自同态，其中该空间定义在代数闭域\(\mathbf{K}\)上。对\(\lambda\)（\(u\)的特征值），令\(\mathrm{E}_{\lambda}\)为\(\mathbf{E}\)中被某个\(u - \lambda \mathrm{I}\)的幂所湮灭的向量组成的子空间。证明若\((\operatorname{ad} u)^n v = 0\)，则\(\mathrm{E}_{\lambda}\)在\(v\)下稳定。
  \end{enumerate}
\end{enumerate}

\begin{enumerate}
\def\labelenumi{(\alph{enumi})}
\setcounter{enumi}{1}
\item
  由 (a) 推出，若 g 是 E 上自同态组成的幂零李代数，则 E 是子空间\(F_{j}\)的直和（\(1 \leqslant j \leqslant m\)），这些子空间在 g 下稳定，并且在每个\(F_{j}\)上，每个\(u \in g\)的限制都可写成\(\lambda_{j}(u)\mathbf{I} + u_{j}\)，其中\(\lambda_{j}(u) \in \mathbf{K}\)，且\(u_{j}\)为幂零的。
\item
  若\(\mathbf{K}\)的特征为 2，\(\mathrm{E} = \mathrm{K}^2\)，且\(\mathfrak{g}\)为幂零代数\(\mathfrak{sl}(2,\mathbf{K})\)，证明\(m = 1\)，并且可能有\(\lambda (u + v)\neq \lambda (u) + \lambda (v)\)，其中\(u,v\)是\(\mathfrak{g}\)中的两个元素。（关于本题的后续内容，参见 § 5，习题 12。）
\end{enumerate}

\begin{enumerate}
\def\labelenumi{\arabic{enumi}.}
\setcounter{enumi}{22}
\tightlist
\item
  令\(\mathfrak{g}\)为李 \(p\)-代数，定义在完美域\(\mathbf{K}\)上，且该域的特征为\(p > 0\)。\(\mathfrak{g}\)称为\(p\)-幂幺的，如果对所有\(x \in \mathfrak{g}\)，存在\(m\)使\(x^{p^m} = 0\)。
\end{enumerate}

\begin{enumerate}
\def\labelenumi{(\alph{enumi})}
\item
  证明每个\(p\)-幂幺李 \(p\)-代数都是幂零的。
\item
  设\(\mathfrak{g}\)为幂零的，并令\(\mathfrak{h}\)表示\(p\)-核，取自\(\mathfrak{g}\)的中心（§ 1，习题 23 (a)）。证明\(\mathfrak{g} / \mathfrak{h}\)是\(p\)-幂幺的。
\item
  令\(\mathfrak{g}\)为幂零李 \(p\)-代数，其三个基元素为\(e_1, e_2, e_3\)，满足\([e_1, e_2] = [e_1, e_3] = 0\)、\([e_2, e_3] = e_1\)、\(e_1^p = e_1\)和\(e_2^p = e_3^p = 0\)。则\(\mathfrak{h} = \mathrm{Ke}_1\)，但\(\mathfrak{g}\)不是\(\mathfrak{h}\)与一个\(p\)-幂幺\(p\)-子代数的直和。
\item
  若 g 是 p-幂幺的，证明在 g 的受限包络代数中，由 g 生成的双边理想是幂零的。（对 g 的维数作归纳。）
\end{enumerate}

\begin{enumerate}
\def\labelenumi{\arabic{enumi}.}
\setcounter{enumi}{23}
\item
  假设 K 的特征为 2。证明习题 9 中的幂零李代数\(\mathfrak{g}_4\)不存在 2-映射。
\item
  令\(\mathfrak{g}\)为李 \(p\)-代数。证明\(\mathfrak{g}\)的最大幂零理想是\(p\)-理想（参见 § 1，习题 22）。
\item
  令\(\mathbf{G}\)为一个群，令\((\mathbf{H}_n)_{n\geqslant 1}\)为\(\mathbf{G}\)的子群递减序列；设\(\mathbf{H}_1 = \mathbf{G}\)，并且若记\((x,y) = xyx^{-1}y^{-1}\)，则关系\(x\in \mathrm{H}_i,y\in \mathrm{H}_j\)蕴含\((x,y)\in \mathrm{H}_{i + j}\)。
\end{enumerate}

\begin{enumerate}
\def\labelenumi{(\alph{enumi})}
\item
  令\(G_{i} = H_{i}/H_{i+1}\)；证明\(G_{i}\)为交换的，并且映射\(x, y \mapsto (x, y)\)在取商后定义了从\(G_{i} \times G_{j}\)到\(G_{i+j}\)的 Z-双线性映射。
\item
  令\(\mathrm{gr}(\mathbf{G}) = \sum_{i=1}^{\infty} \mathbf{G}_i\)，并将 (a) 中定义的映射\(\mathbf{G}_i \times \mathbf{G}_j \to \mathbf{G}_{i+j}\)按线性延拓为从\(\mathrm{gr}(\mathbf{G}) \times \mathrm{gr}(\mathbf{G})\)到\(\mathrm{gr}(\mathbf{G})\)的 Z-双线性映射。证明\(\mathrm{gr}(\mathbf{G})\)关于此映射是李 Z-代数（验证 Jacobi 恒等式时，使用下列公式：
\end{enumerate}

\[
((x, y), z ^ {y}) \cdot ((y, z), x ^ {z}) \cdot ((z, x), y ^ {x}) = e \quad x, y, z \text {   in   } G,
\]

其中\(x^{y}\)表示\(yxy^{-1}\)，\(e\)表示 G 的单位元）。

\begin{enumerate}
\def\labelenumi{(\alph{enumi})}
\setcounter{enumi}{2}
\tightlist
\item
  若存在\(n\)使得\(\mathrm{H}_n = \{e\}\)，证明\(\operatorname{gr}(G)\)是幂零李 \(\mathbf{Z}\)-代数。
\end{enumerate}

\begin{enumerate}
\def\labelenumi{\arabic{enumi}.}
\setcounter{enumi}{26}
\tightlist
\item
   令 A 为带单位元 1 的结合代数，并令\(\mathrm{A}_0 = \mathrm{A} \supset \mathrm{A}_1 \supset \dots \supset \mathrm{A}_n \supset \dots\)为 A 的双边理想递减序列，满足\(\mathrm{A}_i \cdot \mathrm{A}_j \subset \mathrm{A}_{i+j}\)。令 G 为单位元\(e\)的群，并令\(f: \mathrm{G} \to \mathrm{A}\)为映射，满足\(f(e) = 1\)、\(f(xy) = f(x) \cdot f(y)\)，且\(1 - f(x) \in \mathrm{A}_1\)对所有\(x \in \mathrm{G}\)成立。令\(\mathrm{H}_n\)表示所有\(x \in \mathrm{G}\)中满足\(1 - f(x) \in \mathrm{A}_n\)者的集合。证明\(\mathrm{H}_n\)满足习题 26 的条件。证明映射\(x \mapsto f(x) - 1\)在取商后定义了李代数\(\mathrm{gr}(\mathrm{G})\)到与分次环\(\mathrm{gr}(\mathrm{A}) = \sum \mathrm{A}_n / \mathrm{A}_{n+1}\)相关的李代数中的单射同态（参见《交换代数》第三章）。
\end{enumerate}

